\documentclass[11pt]{article}

\usepackage{amsmath, amssymb, amsfonts}
\usepackage{geometry}
\geometry{margin=1in}

\title{Supplementary Information: Posterior Information Distortion, Effective Sample Size, Corrected Covariance, and Bias}
\author{}
\date{}

\begin{document}

\maketitle

% NOTE: This document is the Bayesian/posterior counterpart of the
% Likelihood Information Distortion supplement in
%   theory/macroir/docs/Likelihood_Information_Distortion/
% Sections 1-4 establish definitions for log-posterior, posterior Hessian
% decomposition, posterior score fluctuation matrices, and the Posterior
% Information Distortion (PID). Sections 5+ (sample/correlation distortion,
% distortion-corrected covariance, distortion-induced bias, decomposition,
% summary) are drafted in subsequent commits.

\section*{1. Log-Posterior and Score Definitions}

Consider a trajectory of $T$ sequential measurement intervals with data $D = (y_1,\ldots,y_T)$, and a Gaussian prior $\pi(\theta)$ on the parameter vector $\theta \in \mathbb{R}^d$. Throughout this supplement, all derivatives, scores, and Hessians are taken in the transformed coordinates in which $\pi$ is Gaussian; no Jacobian corrections appear. We restrict to Gaussian priors with fixed (non-random) hyperparameters, so the prior log-density and its derivatives are deterministic functions of $\theta$ only. As in \texttt{supplement\_information\_distortion\_main.tex}, all expectations and covariances are taken under the frequentist sampling distribution of $D$ at a fixed true $\theta$ (the role of the prior here is to regularize curvature and shift the score mean, not to randomize $\theta$).

\paragraph{Notation bridge to the Likelihood supplement.}
We adopt the subscripts $\mathrm{lik}$, $\mathrm{prior}$, $\mathrm{post}$ to distinguish the data, prior, and combined contributions. The likelihood-side objects are renamings of the unsubscripted objects of the Likelihood supplement:
\begin{equation}
\mathbf H_{\mathrm{lik}} \equiv \mathbf H,
\qquad
\mathbf{JT}_{\mathrm{lik}} \equiv \mathbf J_{\mathrm{total}},
\qquad
\mathbf{JS}_{\mathrm{lik}} \equiv \mathbf J_{\mathrm{sample}},
\qquad
S_{\mathrm{lik}} \equiv S,
\label{eq:notation-bridge}
\end{equation}
all of Sections~1--3 of \texttt{supplement\_information\_distortion\_main.tex}. We also write $\ell(\theta) = \log L(\theta)$ as in the Likelihood supplement.

\paragraph{Log-posterior and posterior score.}
The log-posterior decomposes as
\begin{equation}
\log p(\theta \mid D)
=
\ell(\theta) + \log \pi(\theta) + \mathrm{const},
\label{eq:log-posterior-def}
\end{equation}
where $\ell(\theta) = \sum_{t=1}^T \ell_t(\theta)$ with $\ell_t(\theta) = \log p(y_t \mid \theta, y_{1:t-1})$, and the additive constant $-\log p(D)$ does not depend on $\theta$. The posterior score is
\begin{equation}
S_{\mathrm{post}}(\theta)
=
\nabla_\theta \log p(\theta \mid D)
=
S_{\mathrm{lik}}(\theta) + s_\pi(\theta),
\label{eq:post-score-def}
\end{equation}
where
\begin{equation}
S_{\mathrm{lik}}(\theta) = \sum_{t=1}^T s_t(\theta),
\qquad
s_t(\theta) := \nabla_\theta \ell_t(\theta),
\qquad
s_\pi(\theta) := \nabla_\theta \log \pi(\theta).
\label{eq:scores-components}
\end{equation}
For the Gaussian prior $\pi(\theta) = \mathcal{N}(\mu_{\mathrm{prior}}, \Sigma_{\mathrm{prior}})$ with precision $\mathbf H_{\mathrm{prior}} = \Sigma_{\mathrm{prior}}^{-1}$, this gives the explicit form
\begin{equation}
s_\pi(\theta) = -\mathbf H_{\mathrm{prior}}(\theta - \mu_{\mathrm{prior}}).
\label{eq:prior-score-gaussian}
\end{equation}

\paragraph{Per-interval vs.\ global decomposition.}
The likelihood score retains the per-interval structure inherited from the trajectory model, while the prior contributes a single global term that is not indexed by measurement interval. We write
\begin{equation}
S_{\mathrm{post}}(\theta)
=
\Bigl(\sum_{t=1}^T s_t(\theta)\Bigr) + s_\pi(\theta),
\label{eq:post-score-decomp}
\end{equation}
distinguishing the per-interval sum from the global prior term.

\paragraph{Prior as a global deterministic observation.}
Throughout this supplement we treat the prior $\pi(\theta)$ as a single global ``observation'' contributing one deterministic score $s_\pi(\theta)$ and one deterministic curvature $\mathbf H_{\mathrm{prior}} = -\nabla_\theta^2 \log \pi(\theta)$. This convention has the following consequences:
\begin{itemize}
\item At fixed $\theta$, $s_\pi(\theta)$ is a constant under the data-sampling distribution; hence its sampling variance vanishes and its sampling covariance with any $s_t(\theta)$ is zero.
\item The prior contributes additively to the curvature budget without introducing cross-interval terms.
\item For a Gaussian prior, $\mathbf H_{\mathrm{prior}}$ is constant in $\theta$ and diagonal in the transformed coordinates.
\item When $\mathbf H_{\mathrm{prior}}$ is positive definite (proper Gaussian prior) and $\mathbf H_{\mathrm{lik}} \succeq 0$, the posterior curvature $\mathbf H_{\mathrm{post}} = \mathbf H_{\mathrm{lik}} + \mathbf H_{\mathrm{prior}}$ developed in Section~2 is strictly positive definite, so the inverse square roots used in Sections~3--4 are well-defined without subspace projection (cf.\ Section~8 of the Likelihood supplement).
\end{itemize}

\paragraph{Posterior score mean and covariance under sampling.}
Combining \eqref{eq:post-score-def} with the deterministic-prior-score convention,
\begin{align}
\mathbb{E}[S_{\mathrm{post}}(\theta)]
&=
\mathbb{E}[S_{\mathrm{lik}}(\theta)] + s_\pi(\theta),
\label{eq:post-score-mean}\\
\mathrm{Cov}(S_{\mathrm{post}}(\theta))
&=
\mathrm{Cov}(S_{\mathrm{lik}}(\theta))
=
\mathbf{JT}_{\mathrm{lik}},
\label{eq:post-score-cov}
\end{align}
i.e.\ the prior shifts the mean of the posterior score but contributes nothing to its sampling covariance. The covariance contribution of the prior enters the diagnostic framework only through its curvature $\mathbf H_{\mathrm{prior}}$, as defined in Section~3.

\paragraph{Important: $\mathbf{JT}_{\mathrm{post}}$ is NOT $\mathrm{Cov}(S_{\mathrm{post}})$.}
Section~3 will define $\mathbf{JT}_{\mathrm{post}} := \mathbf{JT}_{\mathrm{lik}} + \mathbf H_{\mathrm{prior}}$ by treating the prior as a virtual observation whose deterministic score contributes information $\mathbf H_{\mathrm{prior}}$ to the curvature/fluctuation budget. By \eqref{eq:post-score-cov}, this is \emph{not} the literal covariance of $S_{\mathrm{post}}$: $\mathbf{JT}_{\mathrm{post}} \neq \mathrm{Cov}(S_{\mathrm{post}})$. This definitional choice is what makes the curvature/fluctuation identity $\mathbf{JT}_{\mathrm{post}} = \mathbf H_{\mathrm{post}}$ hold under correct likelihood specification regardless of prior strength.

\paragraph{Reduction to the likelihood case.}
When the prior becomes diffuse $\mathbf H_{\mathrm{prior}} \to 0$ (formally an improper uniform prior, with $s_\pi \to 0$ as an automatic consequence of \eqref{eq:prior-score-gaussian} at any fixed $\theta$), equations \eqref{eq:post-score-def}--\eqref{eq:post-score-cov} collapse to the corresponding likelihood definitions and $S_{\mathrm{post}} \to S_{\mathrm{lik}}$. When $\mathbf H_{\mathrm{lik}}$ is full rank, the limit is regular and every Posterior IDM quantity defined in subsequent sections reduces smoothly to its Likelihood counterpart. When $\mathbf H_{\mathrm{lik}}$ is singular, the limit recovers the subspace-projection regime of Section~8 of the Likelihood supplement; the regularizing role of a strictly positive $\mathbf H_{\mathrm{prior}}$ is exactly what Posterior IDM exploits to avoid that projection.

\section*{2. Posterior Hessian Decomposition}

The negative log-posterior decomposes as
\begin{equation}
-\log p(\theta \mid D)
=
-\ell(\theta)
-\log \pi(\theta)
+ \mathrm{const},
\label{eq:neg-log-posterior}
\end{equation}
so its expected curvature splits additively,
\begin{equation}
\mathbf H_{\mathrm{post}}(\theta)
=
\mathbf H_{\mathrm{lik}}(\theta)
+ \mathbf H_{\mathrm{prior}}(\theta).
\label{eq:posterior-hessian-decomp}
\end{equation}

The likelihood term $\mathbf H_{\mathrm{lik}}$ is the Gaussian Fisher Information of Section~2 of the Likelihood supplement:
\begin{equation}
\mathbf H_{\mathrm{lik}}(\theta)
=
\sum_{t=1}^T \mathbb{E}\!\left[\mathbf I_t^{\mathrm{G}}(\theta)\right]
=
- \mathbb{E}\!\left[\nabla_\theta^2 \ell(\theta)\right],
\label{eq:hlik-def}
\end{equation}
where the second equality is the Fisher identity, valid under correct likelihood specification and standard regularity (interchange of differentiation and expectation). The diagnostic role of $\mathbf H_{\mathrm{lik}}$ versus $\mathbf{JT}_{\mathrm{lik}}$ (Section~3 below) is to quantify departures from this identity.

The prior term is the precision matrix of the Gaussian prior in the transformed coordinates,
\begin{equation}
\mathbf H_{\mathrm{prior}}(\theta)
=
- \nabla_\theta^2 \log \pi(\theta)
= \Sigma_{\mathrm{prior}}^{-1},
\label{eq:hprior-def}
\end{equation}
which is constant in $\theta$. When the prior factorizes across the transformed coordinates, $\mathbf H_{\mathrm{prior}}$ is diagonal with entries $1/\sigma_{\mathrm{prior},i}^2$.

The matrix $\mathbf H_{\mathrm{post}}$ defines the local quadratic geometry of the posterior at the MAP and reduces to $\mathbf H_{\mathrm{lik}}$ in the flat-prior limit $\mathbf H_{\mathrm{prior}} \to 0$ (Section~2.3).

\subsection*{2.1 Positive-Definiteness of $\mathbf H_{\mathrm{post}}$}

A proper Gaussian prior has $\mathbf H_{\mathrm{prior}} \succ 0$. When the analytic Gaussian Fisher Information \eqref{eq:hlik-def} is used, $\mathbf H_{\mathrm{lik}} \succeq 0$ by construction (as a sum of expected outer-product terms). Then for any nonzero $v \in \mathbb{R}^d$,
\begin{equation}
v^\top \mathbf H_{\mathrm{post}}\, v
\;\geq\;
v^\top \mathbf H_{\mathrm{prior}}\, v
\;\geq\;
\lambda_{\min}(\mathbf H_{\mathrm{prior}})\,\|v\|^2
\;>\;0,
\qquad
\mathbf H_{\mathrm{post}}
\;\succeq\;
\mathbf H_{\mathrm{prior}}
\;\succ\; 0,
\label{eq:posterior-hessian-pd}
\end{equation}
so $\mathbf H_{\mathrm{post}}$ is strictly positive definite even when $\mathbf H_{\mathrm{lik}}$ is singular or rank-deficient. No subspace projection is required.

\paragraph{Indefinite estimator-based $\mathbf H_{\mathrm{lik}}$.}
In practice $\mathbf H_{\mathrm{lik}}$ is often approximated by a finite-difference Hessian whose estimator bias can produce negative eigenvalues along weakly identified directions. In that case the Loewner bound $\mathbf H_{\mathrm{post}} \succeq \mathbf H_{\mathrm{prior}}$ no longer follows; instead, positive definiteness of $\mathbf H_{\mathrm{post}}$ requires the prior to dominate the negative spectrum of $\mathbf H_{\mathrm{lik}}$.

Let the spectral decomposition of $\mathbf H_{\mathrm{lik}}$ be
\begin{equation}
\mathbf H_{\mathrm{lik}} = \mathbf H_{\mathrm{lik}}^+ - \mathbf H_{\mathrm{lik}}^-,
\qquad
\mathbf H_{\mathrm{lik}}^+ \succeq \mathbf 0,
\qquad
\mathbf H_{\mathrm{lik}}^- \succeq \mathbf 0,
\label{eq:hlik-jordan-decomposition}
\end{equation}
where $\mathbf H_{\mathrm{lik}}^+$ collects the contribution from non-negative eigenvalues and $\mathbf H_{\mathrm{lik}}^-$ the (sign-flipped) contribution from negative eigenvalues, so that both parts are PSD by construction (Jordan decomposition into positive and negative parts of a symmetric matrix). The necessary-and-sufficient condition for $\mathbf H_{\mathrm{post}} = \mathbf H_{\mathrm{lik}} + \mathbf H_{\mathrm{prior}}$ to be positive definite is
\begin{equation}
\mathbf H_{\mathrm{prior}} - \mathbf H_{\mathrm{lik}}^- \;\succ\; \mathbf 0,
\label{eq:dominance-condition}
\end{equation}
i.e., the prior precision must dominate the negative part of $\mathbf H_{\mathrm{lik}}$ in the Loewner order. Equivalently, for every unit vector $v \in \mathbb{R}^d$,
\begin{equation}
v^\top \mathbf H_{\mathrm{prior}}\, v \;>\; v^\top \mathbf H_{\mathrm{lik}}^-\, v,
\label{eq:dominance-condition-vector}
\end{equation}
i.e., the prior precision in every direction $v$ exceeds the negative part of $\mathbf H_{\mathrm{lik}}$ in that same direction.

\paragraph{Operational sufficient condition.}
A simple sufficient condition that requires only two scalar eigenvalue computations is
\begin{equation}
\lambda_{\min}(\mathbf H_{\mathrm{prior}}) \;>\; \lambda_{\max}(\mathbf H_{\mathrm{lik}}^-)
\;=\;
\max\bigl(0,\, -\lambda_{\min}(\mathbf H_{\mathrm{lik}})\bigr),
\label{eq:dominance-scalar-bound}
\end{equation}
which is sharp (necessary and sufficient) when $\mathbf H_{\mathrm{prior}}$ and $\mathbf H_{\mathrm{lik}}$ are simultaneously diagonalizable, and otherwise a conservative bound. In the codebase, the prior is diagonal in the transformed coordinates while $\mathbf H_{\mathrm{lik}}$ generally has off-diagonal entries, so \eqref{eq:dominance-scalar-bound} is conservative and \eqref{eq:dominance-condition} is the diagnostic of record. The sufficient bound \eqref{eq:dominance-scalar-bound} is convenient as a quick monitoring statistic, but a rejection by \eqref{eq:dominance-scalar-bound} must trigger the matrix check \eqref{eq:dominance-condition} (cheap: a single eigendecomposition of $\mathbf H_{\mathrm{lik}}$ plus the eigenvalue test of $\mathbf H_{\mathrm{prior}} - \mathbf H_{\mathrm{lik}}^-$) before declaring $\mathbf H_{\mathrm{post}}$ ill-defined.

\subsection*{2.2 Consequences for Posterior IDM}

Equation \eqref{eq:posterior-hessian-pd} has two structural consequences that distinguish Posterior IDM from its Likelihood counterpart:
\begin{itemize}
\item All Posterior IDM quantities defined below ($\mathbf C_{\mathrm{post}}$, $\mathbf C_{\mathrm{sample},\mathrm{post}}$, $\mathbf R_{\mathrm{post}}$, $\Sigma_{\mathrm{DCC},\mathrm{post}}$, $\mathbf b_{\mathrm{DIB},\mathrm{post}}$) are well-defined on the full parameter space whenever $\mathbf H_{\mathrm{lik}} \succeq 0$ (analytic Gaussian Fisher case) or whenever \eqref{eq:dominance-condition} holds (finite-difference case). The subspace projection of Section~8 of the Likelihood supplement is unnecessary in either case.
\item The whitening transform is computed via the spectral decomposition $\mathbf H_{\mathrm{post}} = \mathbf U_{\mathrm{post}}\, \boldsymbol\Lambda_{\mathrm{post}}\, \mathbf U_{\mathrm{post}}^\top$, yielding the symmetric inverse square root $\mathbf H_{\mathrm{post}}^{-1/2} = \mathbf U_{\mathrm{post}}\, \boldsymbol\Lambda_{\mathrm{post}}^{-1/2}\, \mathbf U_{\mathrm{post}}^\top$. Both per-eigen-direction (eigenvalues, $\mathrm{tr}$, $\det$) and per-parameter ($[\mathbf C_{\mathrm{post}}]_{ii}$) diagnostics are then basis-invariant and interpretable as in Section~5 of the Likelihood supplement; see Section~4.5 for the operational details.
\end{itemize}

In computations we never form $\mathbf H_{\mathrm{lik}}$ and $\mathbf H_{\mathrm{prior}}$ separately when a stable factor of $\mathbf H_{\mathrm{post}}$ is required: we assemble the sum \eqref{eq:posterior-hessian-decomp} first, then factor.

\subsection*{2.3 Limiting Regimes}

\paragraph{Flat-prior limit.}
As $\mathbf H_{\mathrm{prior}} \to 0$, \eqref{eq:posterior-hessian-decomp} degenerates to $\mathbf H_{\mathrm{post}} \to \mathbf H_{\mathrm{lik}}$. Combined with $\mathbf{JT}_{\mathrm{post}} \to \mathbf{JT}_{\mathrm{lik}}$ and $\mathbf{JS}_{\mathrm{post}} \to \mathbf{JS}_{\mathrm{lik}}$ from Section~3, every Posterior IDM quantity reduces termwise to its Likelihood IDM counterpart. When $\mathbf H_{\mathrm{lik}}$ is rank-deficient, this collapse is understood on the retained eigenspace of $\mathbf H_{\mathrm{lik}}$ (the singular-Fisher machinery of Section~8 of the Likelihood supplement re-emerges as the appropriate generalization).

\paragraph{Informative-prior limit.}
In the opposite limit $\lambda_{\min}(\mathbf H_{\mathrm{prior}}) \to \infty$, $\mathbf H_{\mathrm{post}}$ is dominated by $\mathbf H_{\mathrm{prior}}$, the posterior contracts to the prior, and the conditioning of $\mathbf H_{\mathrm{post}}$ is set by the prior precision. All Posterior IDM quantities remain well-defined throughout this entire range.

\section*{3. Posterior Score Fluctuation Matrices}

The diagnostic role played by $\mathbf{JT}_{\mathrm{lik}}$ ($\equiv \mathbf J_{\mathrm{total}}$) and $\mathbf{JS}_{\mathrm{lik}}$ ($\equiv \mathbf J_{\mathrm{sample}}$) for the likelihood (Section~3 of the Likelihood supplement) must be lifted to the posterior so that the comparison with $\mathbf H_{\mathrm{post}}$ yields meaningful values.

Recall the posterior score \eqref{eq:post-score-def}. Taken over realizations of $D$ at fixed $\theta$, $s_\pi(\theta)$ is a constant; hence it has zero sampling variance and, as a consequence, zero sampling covariance with $S_{\mathrm{lik}}(\theta)$. The raw posterior score covariance therefore satisfies
\begin{equation}
\mathrm{Cov}\!\left(S_{\mathrm{post}}(\theta)\right)
=
\mathrm{Cov}\!\left(S_{\mathrm{lik}}(\theta)\right)
=
\mathbf{JT}_{\mathrm{lik}},
\label{eq:posterior-score-cov-raw}
\end{equation}
i.e.\ it does \emph{not} contain $\mathbf H_{\mathrm{prior}}$, while the negative posterior Hessian $\mathbf H_{\mathrm{post}} = \mathbf H_{\mathrm{lik}} + \mathbf H_{\mathrm{prior}}$ does. A direct comparison $\mathbf H_{\mathrm{post}}^{-1/2} \mathbf{JT}_{\mathrm{lik}} \mathbf H_{\mathrm{post}}^{-1/2}$ would conflate prior shrinkage with genuine information distortion and would not equal $\mathbf I$ under correct likelihood specification.

\subsection*{3.1 Prior as a single deterministic observation}

We adopt the convention of treating the prior as one additional observation, independent of the data, with precision $\mathbf H_{\mathrm{prior}}$ and a score that is deterministic across data realizations (zero fluctuation in $D$ at fixed $\theta$). Under this convention, the prior contributes its precision to every information-style matrix but contributes no temporal cross-terms with the data, since it is fixed before $D$ is observed. We accordingly define
\begin{equation}
\mathbf{JT}_{\mathrm{post}}
\;:=\;
\mathbf{JT}_{\mathrm{lik}} + \mathbf H_{\mathrm{prior}},
\label{eq:JT-post-def}
\end{equation}
\begin{equation}
\mathbf{JS}_{\mathrm{post}}
\;:=\;
\mathbf{JS}_{\mathrm{lik}} + \mathbf H_{\mathrm{prior}}.
\label{eq:JS-post-def}
\end{equation}

\paragraph{$\mathbf{JT}_{\mathrm{post}}$ is a defined object, not a covariance.}
By \eqref{eq:posterior-score-cov-raw}, $\mathbf{JT}_{\mathrm{post}}$ is \emph{not} equal to the literal covariance of $S_{\mathrm{post}}(\theta)$ (which is $\mathbf{JT}_{\mathrm{lik}}$). Definition \eqref{eq:JT-post-def} is the augmented-observation bookkeeping convention that places the prior on the same footing as the data: it contributes its curvature directly to the information budget rather than through sampling variance. This convention is what propagates the likelihood Bartlett identity to the posterior in \eqref{eq:posterior-info-identity} below.

\paragraph{Inheritance of the information identity.}
With \eqref{eq:JT-post-def}, the posterior information identity inherits from the likelihood one. If $\mathbf{JT}_{\mathrm{lik}} = \mathbf H_{\mathrm{lik}}$ (correct likelihood specification, Section~4 of the Likelihood supplement), then
\begin{equation}
\mathbf{JT}_{\mathrm{post}}
=
\mathbf H_{\mathrm{lik}} + \mathbf H_{\mathrm{prior}}
=
\mathbf H_{\mathrm{post}}.
\label{eq:posterior-info-identity}
\end{equation}
This identity holds by construction of $\mathbf{JT}_{\mathrm{post}}$ (Eq.~\eqref{eq:JT-post-def}), not as an independent Bartlett identity for the posterior score; it is the additive convention that propagates the likelihood identity. Crucially, the equality survives arbitrary rescaling of $\mathbf H_{\mathrm{prior}}$: as the prior ranges from weakly to dominantly informative, $\mathbf{JT}_{\mathrm{post}}$ and $\mathbf H_{\mathrm{post}}$ change in lockstep, so the resulting $\mathbf C_{\mathrm{post}} = \mathbf I$ in Section~4 holds regardless of prior strength.

\subsection*{3.2 Sample/total decomposition}

There is no posterior per-interval score $s_{t,\mathrm{post}}$; the prior is treated as a single extra observation, so the within-interval covariance sum in $\mathbf{JS}_{\mathrm{post}}$ is over data intervals only. Since the prior pseudo-score is deterministic, $\mathrm{Cov}(s_t, s_\pi) = 0$ for every data interval $t$ and the cross-interval sum runs only over data indices. The decomposition of $\mathbf{JT}_{\mathrm{lik}}$ into within-interval variance plus cross-interval covariance therefore lifts unchanged:
\begin{equation}
\mathbf{JT}_{\mathrm{post}}
=
\mathbf{JS}_{\mathrm{post}}
+
2\sum_{i<j}\mathrm{Cov}(s_i, s_j),
\qquad i,j \in \{1,\ldots,T\},
\label{eq:posterior-sample-total}
\end{equation}
with $s_t$ the per-interval data score of \eqref{eq:scores-components}. The prior contributes $\mathbf H_{\mathrm{prior}}$ identically to both $\mathbf{JT}_{\mathrm{post}}$ and $\mathbf{JS}_{\mathrm{post}}$, and therefore drops out of the difference $\mathbf{JT}_{\mathrm{post}} - \mathbf{JS}_{\mathrm{post}}$, which equals $2\sum_{i<j}\mathrm{Cov}(s_i, s_j)$ and involves only the data scores. Equation \eqref{eq:posterior-sample-total} isolates temporal dependence in the data score from the prior contribution, exactly as in the likelihood case.

\subsection*{3.3 Vanishing prior limit}

As $\mathbf H_{\mathrm{prior}} \to \mathbf 0$,
\begin{equation}
\mathbf{JT}_{\mathrm{post}} \to \mathbf{JT}_{\mathrm{lik}},
\qquad
\mathbf{JS}_{\mathrm{post}} \to \mathbf{JS}_{\mathrm{lik}},
\qquad
\mathbf H_{\mathrm{post}} \to \mathbf H_{\mathrm{lik}},
\label{eq:flat-prior-limit}
\end{equation}
and all posterior distortion matrices defined in Section~4 reduce to their Likelihood IDM counterparts. When $\mathbf H_{\mathrm{lik}}$ is full rank the limit is direct; when $\mathbf H_{\mathrm{lik}}$ is rank-deficient, the limit must be interpreted on its retained eigenspace (Section~8 of the Likelihood supplement). For strictly positive $\mathbf H_{\mathrm{prior}}$, by contrast, the Posterior IDM is well-defined globally without any retained-subspace construction; this is the structural advantage exploited in Section~4.

\section*{4. Posterior Information Distortion (PID)}

The central Bayesian diagnostic is the total Posterior Information Distortion
\begin{equation}
\mathbf C_{\mathrm{post}}(\theta)
=
\mathbf H_{\mathrm{post}}^{-1/2}
\mathbf{JT}_{\mathrm{post}}
\mathbf H_{\mathrm{post}}^{-1/2},
\label{eq:posterior-c-def}
\end{equation}
the direct counterpart of the Likelihood Information Distortion
$\mathbf C = \mathbf H^{-1/2} \mathbf J_{\mathrm{total}} \mathbf H^{-1/2}$
(Section~4 of the Likelihood supplement; recall the bridge \eqref{eq:notation-bridge} that identifies $\mathbf H \equiv \mathbf H_{\mathrm{lik}}$ and $\mathbf J_{\mathrm{total}} \equiv \mathbf{JT}_{\mathrm{lik}}$), with the likelihood curvature and total score covariance replaced by their posterior counterparts
\begin{equation}
\mathbf H_{\mathrm{post}} = \mathbf H_{\mathrm{lik}} + \mathbf H_{\mathrm{prior}},
\qquad
\mathbf{JT}_{\mathrm{post}} = \mathbf{JT}_{\mathrm{lik}} + \mathbf H_{\mathrm{prior}}.
\label{eq:posterior-h-j}
\end{equation}
Throughout, $\mathbf H_{\mathrm{lik}}$, $\mathbf H_{\mathrm{prior}}$, $\mathbf{JT}_{\mathrm{lik}}$, and $\mathbf{JT}_{\mathrm{post}}$ are symmetric, and the symmetric inverse square root $\mathbf H_{\mathrm{post}}^{-1/2}$ is the unique positive definite root of $\mathbf H_{\mathrm{post}}^{-1}$.

\subsection*{4.1 Consistency under correct likelihood specification}

Under correct likelihood specification and standard regularity at the true parameter $\theta_0$, the population information identity $\mathbf{JT}_{\mathrm{lik}} = \mathbf H_{\mathrm{lik}}$ holds; empirical plug-ins satisfy it only up to sampling fluctuation. Substituting into \eqref{eq:posterior-h-j} gives \eqref{eq:posterior-info-identity}, $\mathbf{JT}_{\mathrm{post}} = \mathbf H_{\mathrm{post}}$, and therefore
\begin{equation}
\mathbf C_{\mathrm{post}}
=
\mathbf H_{\mathrm{post}}^{-1/2}
\mathbf H_{\mathrm{post}}
\mathbf H_{\mathrm{post}}^{-1/2}
=
\mathbf I.
\label{eq:posterior-c-identity}
\end{equation}
This holds for \emph{any} proper Gaussian prior, informative or weak: the prior contribution $\mathbf H_{\mathrm{prior}}$ enters $\mathbf H_{\mathrm{post}}$ and $\mathbf{JT}_{\mathrm{post}}$ symmetrically and cancels under the whitening transform.

The companion diagnostics $\mathbf C_{\mathrm{sample},\mathrm{post}} = \mathbf H_{\mathrm{post}}^{-1/2} \mathbf{JS}_{\mathrm{post}} \mathbf H_{\mathrm{post}}^{-1/2}$ and $\mathbf R_{\mathrm{post}} = \mathbf{JS}_{\mathrm{post}}^{-1/2} \mathbf{JT}_{\mathrm{post}} \mathbf{JS}_{\mathrm{post}}^{-1/2}$ play the same role for within-interval mismatch and temporal correlation as $\mathbf C_{\mathrm{sample}}$ and $\mathbf R$ do in the Likelihood supplement. Under correct specification with no temporal correlation, $\mathbf C_{\mathrm{post}} = \mathbf C_{\mathrm{sample},\mathrm{post}} = \mathbf R_{\mathrm{post}} = \mathbf I$; under correct specification with temporal correlation, $\mathbf C_{\mathrm{sample},\mathrm{post}} = \mathbf I$ and $\mathbf R_{\mathrm{post}} \neq \mathbf I$, with $\mathbf C_{\mathrm{post}}$ tracking $\mathbf R_{\mathrm{post}}$ (the formal development is deferred to the sample/correlation decomposition section below).

\subsection*{4.2 Deviation as a pure likelihood diagnostic}

If the likelihood is misspecified, $\mathbf{JT}_{\mathrm{lik}} \neq \mathbf H_{\mathrm{lik}}$ and
\begin{equation}
\mathbf{JT}_{\mathrm{post}} - \mathbf H_{\mathrm{post}}
=
\mathbf{JT}_{\mathrm{lik}} - \mathbf H_{\mathrm{lik}},
\label{eq:posterior-residual}
\end{equation}
so the prior contribution drops out of the unwhitened residual. Equivalently,
\begin{equation}
\mathbf C_{\mathrm{post}} - \mathbf I
=
\mathbf H_{\mathrm{post}}^{-1/2}
\bigl(\mathbf{JT}_{\mathrm{lik}} - \mathbf H_{\mathrm{lik}}\bigr)
\mathbf H_{\mathrm{post}}^{-1/2}.
\label{eq:posterior-c-residual}
\end{equation}
The numerator $\mathbf{JT}_{\mathrm{lik}} - \mathbf H_{\mathrm{lik}}$ depends only on likelihood quantities, so any nonzero deviation $\mathbf C_{\mathrm{post}} \neq \mathbf I$ originates purely on the likelihood side: the prior cannot create distortion of this kind. The whitening matrix $\mathbf H_{\mathrm{post}}^{-1/2}$ does, however, rescale the magnitude of the deviation: a strong prior shrinks $\mathbf C_{\mathrm{post}} - \mathbf I$ toward zero (in the limit $\lambda_{\min}(\mathbf H_{\mathrm{prior}}) \to \infty$, $\mathbf C_{\mathrm{post}} \to \mathbf I$), while a weak prior recovers the Likelihood IDM magnitude $\mathbf C - \mathbf I = \mathbf H_{\mathrm{lik}}^{-1/2}(\mathbf{JT}_{\mathrm{lik}} - \mathbf H_{\mathrm{lik}})\mathbf H_{\mathrm{lik}}^{-1/2}$ (Section~4.3). Thus $\mathbf C_{\mathrm{post}}$ diagnoses purely likelihood-side miscalibration, with a magnitude calibrated by total posterior curvature.

\subsection*{4.3 Flat-prior limit}

As $\mathbf H_{\mathrm{prior}} \to \mathbf 0$,
\begin{equation}
\mathbf H_{\mathrm{post}} \to \mathbf H_{\mathrm{lik}},
\qquad
\mathbf{JT}_{\mathrm{post}} \to \mathbf{JT}_{\mathrm{lik}},
\qquad
\mathbf C_{\mathrm{post}} \to \mathbf C,
\label{eq:posterior-flat-limit}
\end{equation}
recovering the Likelihood IDM of Section~4 of the Likelihood supplement. When $\mathbf H_{\mathrm{lik}}$ is full rank this limit is literal; when $\mathbf H_{\mathrm{lik}}$ is rank-deficient it is understood on its retained eigenspace (Section~8 of the Likelihood supplement), and the Posterior IDM with a vanishing-but-proper prior provides a regularized version of that subspace construction. The Posterior IDM is therefore a strict generalization that reduces to its frequentist analogue in the noninformative limit.

\subsection*{4.4 Well-posedness without subspace projection}

For any proper Gaussian prior, $\mathbf H_{\mathrm{prior}} \succ 0$ (diagonal in transformed coordinates). Combined with $\mathbf H_{\mathrm{lik}} \succeq 0$ for the analytic Gaussian Fisher Information, Eq.\ \eqref{eq:posterior-hessian-pd} gives
\begin{equation}
\mathbf H_{\mathrm{post}} \succeq \mathbf H_{\mathrm{prior}} \succ \mathbf 0,
\qquad \text{so}\qquad
\mathbf H_{\mathrm{post}}^{-1/2}\ \text{in \eqref{eq:posterior-c-def} is unconditionally defined.}
\label{eq:posterior-c-wellposed}
\end{equation}
Hence $\mathbf C_{\mathrm{post}}$ requires no retained-eigenspace construction. When $\mathbf H_{\mathrm{lik}}$ is a finite-difference estimator that can develop negative eigenvalues, well-posedness instead requires the explicit dominance condition \eqref{eq:dominance-condition}; for proper priors and bounded estimator error this is easily monitored. This contrasts with the Likelihood IDM, where singular $\mathbf H_{\mathrm{lik}}$ forces the subspace treatment of Section~8 of the Likelihood supplement.

\paragraph{Numerical caveat.}
Mathematical well-posedness does not imply numerical robustness: directions along which both $\mathbf H_{\mathrm{lik}} \approx 0$ and the prior precision is small produce $\mathbf H_{\mathrm{post}}$ with large condition number, and $\mathbf C_{\mathrm{post}}$ entries in those directions are sensitive to prior specification. We recommend reporting the condition number of $\mathbf H_{\mathrm{post}}$ alongside any $\mathbf C_{\mathrm{post}}$ diagnostic.

\subsection*{4.5 Computation}

We never form $\mathbf H_{\mathrm{post}}^{-1}$ explicitly. The symmetric inverse square root is obtained from the spectral decomposition
\begin{equation}
\mathbf H_{\mathrm{post}} = \mathbf U_{\mathrm{post}}\, \boldsymbol\Lambda_{\mathrm{post}}\, \mathbf U_{\mathrm{post}}^\top,
\qquad
\mathbf H_{\mathrm{post}}^{-1/2}
=
\mathbf U_{\mathrm{post}}\,
\boldsymbol\Lambda_{\mathrm{post}}^{-1/2}\,
\mathbf U_{\mathrm{post}}^\top,
\label{eq:posterior-symm-sqrt}
\end{equation}
where $\mathbf U_{\mathrm{post}}$ is orthogonal and $\boldsymbol\Lambda_{\mathrm{post}}$ is diagonal positive (modulo the dominance condition \eqref{eq:dominance-condition} in the finite-difference case). The whitening transform of \eqref{eq:posterior-c-def} is then computed in retained-eigenspace coordinates: project $\mathbf{JT}_{\mathrm{post}}$ onto $\mathbf U_{\mathrm{post}}$, apply the diagonal scaling $\boldsymbol\Lambda_{\mathrm{post}}^{-1/2}$ on both sides, and embed back through $\mathbf U_{\mathrm{post}}$. This is numerically stable, produces a symmetric result, and degenerates gracefully when $\mathbf H_{\mathrm{lik}}$ contributes a near-singular direction (the prior precision still fixes that direction; no subspace projection required as long as the dominance condition \eqref{eq:dominance-condition} holds).

\paragraph{Per-parameter vs.\ spectral diagnostics.}
Because the symmetric inverse square root \eqref{eq:posterior-symm-sqrt} is used throughout, $\mathbf C_{\mathrm{post}}$ as computed is the SPD matrix of \eqref{eq:posterior-c-def}, not a basis-dependent congruent variant. Both per-parameter diagonals $[\mathbf C_{\mathrm{post}}]_{ii}$ and spectral invariants ($\lambda_i$, $\mathrm{tr}$, $\det$) are therefore well-defined and basis-invariant in the same sense as Section~5 of the Likelihood supplement:
\begin{itemize}
\item $[\mathbf C_{\mathrm{post}}]_{ii} = 1$ indicates posterior consistency along parameter $i$; values $>1$ ($<1$) indicate under- (over-) estimated posterior uncertainty relative to the Bayesian curvature $\mathbf H_{\mathrm{post}}$.
\item $\lambda_i(\mathbf C_{\mathrm{post}}) = 1$ indicates consistency along the $i$-th eigen-direction; deviations have the same inflation/deflation interpretation as in Section~5 of the Likelihood supplement.
\end{itemize}



\section*{5. Per-Parameter and Spectral Distortion}

The Posterior Information Distortion $\mathbf C_{\mathrm{post}}$ of Section~4 admits two complementary readouts: a per-parameter diagonal in the original parameter basis and a spectral profile in its eigenbasis. Both are well-defined and basis-invariant under the symmetric whitening of Section~4.5, in the same sense as Section~5 of the Likelihood supplement, but with the likelihood curvature and total score covariance replaced by their posterior counterparts of Eq.~\eqref{eq:posterior-h-j}.

\subsection*{5.1 Per-Parameter Distortion}

In the original parameter basis, the diagonal entries of $\mathbf C_{\mathrm{post}}$ read as a per-parameter ratio of realized score fluctuation to Bayesian curvature-predicted variance:
\begin{equation}
[\mathbf C_{\mathrm{post}}]_{ii}
\;\sim\;
\frac{\text{real fluctuation variance in direction } i}
     {\text{Gaussian-posterior predicted variance in direction } i},
\label{eq:posterior-c-diag-interp}
\end{equation}
where the predicted variance is the direction-$i$ component of $\mathbf H_{\mathrm{post}}^{-1}$ and the fluctuation variance is the corresponding component of $\mathbf{JT}_{\mathrm{post}}$, both measured in the whitened metric of \eqref{eq:posterior-c-def}. Strictly, $[\mathbf C_{\mathrm{post}}]_{ii} = e_i^\top \mathbf H_{\mathrm{post}}^{-1/2}\, \mathbf{JT}_{\mathrm{post}}\, \mathbf H_{\mathrm{post}}^{-1/2} e_i$, which coincides with the literal coordinate ratio $[\mathbf{JT}_{\mathrm{post}}]_{ii}/[\mathbf H_{\mathrm{post}}]_{ii}$ only when $\mathbf H_{\mathrm{post}}$ is diagonal in the parameter basis; in general \eqref{eq:posterior-c-diag-interp} is the corresponding quantity in the whitened metric. With this reading:
\begin{itemize}
\item $[\mathbf C_{\mathrm{post}}]_{ii} = 1$: posterior consistency for parameter $i$.
\item $[\mathbf C_{\mathrm{post}}]_{ii} > 1$: under-estimation of posterior variance for parameter $i$ (over-confident Gaussian posterior).
\item $[\mathbf C_{\mathrm{post}}]_{ii} < 1$: over-estimation of posterior variance for parameter $i$ (under-confident Gaussian posterior).
\end{itemize}

\paragraph{Basis invariance of the diagonal.}
Because $\mathbf C_{\mathrm{post}}$ is built from the \emph{symmetric} inverse square root $\mathbf H_{\mathrm{post}}^{-1/2}$ of Eq.~\eqref{eq:posterior-symm-sqrt}, an orthogonal change of the original parameter basis $\mathbf Q$ acts congruently, $\mathbf H_{\mathrm{post}} \mapsto \mathbf Q\, \mathbf H_{\mathrm{post}}\, \mathbf Q^\top$ and $\mathbf{JT}_{\mathrm{post}} \mapsto \mathbf Q\, \mathbf{JT}_{\mathrm{post}}\, \mathbf Q^\top$, so $\mathbf C_{\mathrm{post}} \mapsto \mathbf Q\, \mathbf C_{\mathrm{post}}\, \mathbf Q^\top$. The diagonal $[\mathbf C_{\mathrm{post}}]_{ii}$ in any chosen parameter basis and the spectrum of $\mathbf C_{\mathrm{post}}$ are both invariant under this group action. A Cholesky whitening $\mathbf H_{\mathrm{post}} = \mathbf L_{\mathrm{post}}\mathbf L_{\mathrm{post}}^\top$ would instead yield a basis-dependent congruent matrix whose diagonal mixes parameter directions through the lower-triangular structure of $\mathbf L_{\mathrm{post}}$; the symmetric construction avoids this.

\subsection*{5.2 Spectral Distortion}

Let $\lambda_i = \lambda_i(\mathbf C_{\mathrm{post}})$ denote the eigenvalues of $\mathbf C_{\mathrm{post}}$, ordered $\lambda_1 \geq \cdots \geq \lambda_d > 0$. By construction $\mathbf C_{\mathrm{post}}$ is symmetric positive definite (because $\mathbf H_{\mathrm{post}} \succ \mathbf 0$ by Section~4.4 and $\mathbf{JT}_{\mathrm{post}} = \mathbf{JT}_{\mathrm{lik}} + \mathbf H_{\mathrm{prior}} \succ \mathbf 0$ as the sum of a PSD covariance and a positive prior precision; congruence preserves positive definiteness), so the spectrum is real and strictly positive. Each $\lambda_i$ admits the same inflation/deflation reading as in Section~5 of the Likelihood supplement, applied along the corresponding eigen-direction of $\mathbf C_{\mathrm{post}}$:
\begin{itemize}
\item $\lambda_i = 1$: perfect agreement between Bayesian curvature and posterior score fluctuation along eigen-direction $i$.
\item $\lambda_i > 1$: variance inflation along eigen-direction $i$ (the Gaussian posterior under-estimates uncertainty there).
\item $\lambda_i < 1$: deflation along eigen-direction $i$ (the Gaussian posterior over-estimates uncertainty there).
\end{itemize}

\paragraph{Global summaries.}
The determinant and trace of $\mathbf C_{\mathrm{post}}$ collapse the spectrum to scalar diagnostics:
\begin{equation}
\det(\mathbf C_{\mathrm{post}})
=
\prod_{i=1}^d \lambda_i,
\qquad
\frac{1}{d}\,\mathrm{tr}(\mathbf C_{\mathrm{post}})
=
\frac{1}{d}\sum_{i=1}^d \lambda_i.
\label{eq:posterior-c-global}
\end{equation}
The determinant quantifies the overall posterior information-volume distortion: $\det(\mathbf C_{\mathrm{post}}) = 1$ under the consistency identity \eqref{eq:posterior-c-identity}, $\det(\mathbf C_{\mathrm{post}}) > 1$ signals net volumetric under-estimation of posterior uncertainty, and $\det(\mathbf C_{\mathrm{post}}) < 1$ signals net volumetric over-estimation. The average directional distortion $\mathrm{tr}(\mathbf C_{\mathrm{post}})/d$ is the mean per-eigen-direction inflation, useful as a single lightweight summary when the spectrum is concentrated. The geometric-mean summary $\exp\bigl(\tfrac{1}{d}\log\det(\mathbf C_{\mathrm{post}})\bigr)$ is the more robust alternative when the spectrum is heavy-tailed.

\paragraph{Correct-specification readout.}
Under the population information identity $\mathbf{JT}_{\mathrm{lik}} = \mathbf H_{\mathrm{lik}}$ at $\theta_0$, Eq.~\eqref{eq:posterior-c-identity} gives $\mathbf C_{\mathrm{post}} = \mathbf I$, hence
\begin{equation}
[\mathbf C_{\mathrm{post}}]_{ii} = 1\ \forall i,
\qquad
\lambda_i(\mathbf C_{\mathrm{post}}) = 1\ \forall i,
\qquad
\det(\mathbf C_{\mathrm{post}}) = 1,
\qquad
\mathrm{tr}(\mathbf C_{\mathrm{post}})/d = 1,
\label{eq:posterior-c-correct-spec}
\end{equation}
\emph{regardless of prior strength}: every per-parameter and spectral summary collapses to its consistency value.

\subsection*{5.3 Limiting Regimes}

\paragraph{Flat-prior limit.}
As $\mathbf H_{\mathrm{prior}} \to \mathbf 0$, both whitening matrix and numerator in \eqref{eq:posterior-c-def} converge to their Likelihood IDM counterparts (Section~4.3 above), and the per-parameter and spectral diagnostics reduce termwise:
\begin{equation}
[\mathbf C_{\mathrm{post}}]_{ii} \to [\mathbf C]_{ii},
\qquad
\lambda_i(\mathbf C_{\mathrm{post}}) \to \lambda_i(\mathbf C),
\qquad
\det(\mathbf C_{\mathrm{post}}) \to \det(\mathbf C),
\qquad
\mathrm{tr}(\mathbf C_{\mathrm{post}})/d \to \mathrm{tr}(\mathbf C)/d,
\label{eq:posterior-c-flat-limit}
\end{equation}
recovering Section~5 of the Likelihood supplement. Diagonal convergence is immediate from entrywise convergence of $\mathbf C_{\mathrm{post}}$; eigenvalue convergence follows from Weyl's inequality under operator-norm convergence of $\mathbf C_{\mathrm{post}}$ to $\mathbf C$. When $\mathbf H_{\mathrm{lik}}$ is rank-deficient the limit is on its retained eigenspace (Section~8 of the Likelihood supplement); a strictly positive $\mathbf H_{\mathrm{prior}}$ is exactly what makes the per-parameter and spectral diagnostics globally well-defined throughout.

\paragraph{Strong-prior limit.}
Loewner comparisons between $\mathbf H_{\mathrm{prior}}$ and $\mathbf H_{\mathrm{lik}}$ alone do not control $\mathbf C_{\mathrm{post}}$ per direction: the symmetric whitening $\mathbf H_{\mathrm{post}}^{-1/2} (\cdot)\, \mathbf H_{\mathrm{post}}^{-1/2}$ is a non-local operation, and the relevant magnitude is that of $\mathbf H_{\mathrm{post}}$ itself relative to the residual $\mathbf{JT}_{\mathrm{lik}} - \mathbf H_{\mathrm{lik}}$. Quantitatively, the residual identity \eqref{eq:posterior-c-residual} gives
\begin{equation}
\|\mathbf C_{\mathrm{post}} - \mathbf I\|
\;\leq\;
\lambda_{\min}(\mathbf H_{\mathrm{post}})^{-1}\,
\bigl\|\mathbf{JT}_{\mathrm{lik}} - \mathbf H_{\mathrm{lik}}\bigr\|,
\label{eq:posterior-c-residual-bound}
\end{equation}
so as $\lambda_{\min}(\mathbf H_{\mathrm{prior}}) \to \infty$ (with $\mathbf{JT}_{\mathrm{lik}} - \mathbf H_{\mathrm{lik}}$ bounded), $\mathbf C_{\mathrm{post}} \to \mathbf I$ uniformly, hence $[\mathbf C_{\mathrm{post}}]_{ii} \to 1$ and $\lambda_i(\mathbf C_{\mathrm{post}}) \to 1$ for every $i$. The diagnostic reports values close to one in the strong-prior regime, where the likelihood-side mismatch is whitened away by the large $\mathbf H_{\mathrm{post}}$, in agreement with Section~4.2.

\paragraph{Large-data (data-dominated) limit.}
Under $N$ independent trajectories at fixed proper $\mathbf H_{\mathrm{prior}}$, $\mathbf H_{\mathrm{lik}} = O(N)$ and $\mathbf{JT}_{\mathrm{lik}} = O(N)$ while $\mathbf H_{\mathrm{prior}} = O(1)$. Hence $\mathbf H_{\mathrm{post}}/N \to \mathbf H_{\mathrm{lik}}/N$ and $\mathbf{JT}_{\mathrm{post}}/N \to \mathbf{JT}_{\mathrm{lik}}/N$ (with the prior contributions $O(1/N)$), so
\begin{equation}
[\mathbf C_{\mathrm{post}}]_{ii} \to [\mathbf C]_{ii},
\qquad
\lambda_i(\mathbf C_{\mathrm{post}}) \to \lambda_i(\mathbf C),
\qquad
N \to \infty,
\label{eq:posterior-c-large-n}
\end{equation}
the data-dominated counterpart of \eqref{eq:posterior-c-flat-limit}. Together with the correct-specification readout \eqref{eq:posterior-c-correct-spec} and the strong-prior limit \eqref{eq:posterior-c-residual-bound}, the per-parameter and spectral diagnostics behave coherently across all three boundary regimes: $\mathbf C_{\mathrm{post}} = \mathbf I$ under correct specification (any prior), $\mathbf C_{\mathrm{post}} \to \mathbf I$ under strong prior (any likelihood), $\mathbf C_{\mathrm{post}} \to \mathbf C$ under data-dominated or flat-prior asymptotics. The full-rank caveats of Sections~4.3 (analytic Gaussian Fisher) and 4.4 (dominance condition \eqref{eq:dominance-condition} for finite-difference $\mathbf H_{\mathrm{lik}}$) carry over without change.

\section*{6. Posterior Distortion-Corrected Covariance (DCC)}

The Bayesian counterpart of the Likelihood DCC of Section~6 of the Likelihood supplement is obtained by repeating the first-order Taylor expansion argument at the maximum a posteriori (MAP) estimate $\hat\theta_{\mathrm{MAP}}$ rather than at the MLE. As in Section~4, the prior contribution enters both the curvature ($\mathbf H_{\mathrm{post}}$) and the information budget ($\mathbf{JT}_{\mathrm{post}}$); the resulting sandwich form is the Posterior DCC, $\Sigma_{\mathrm{DCC},\mathrm{post}}$ (formally defined in \eqref{eq:posterior-sigma-dcc-def} below).

\subsection*{6.1 Taylor Expansion of the Posterior Score at the MAP}

Expand the posterior score $S_{\mathrm{post}}(\theta) = S_{\mathrm{lik}}(\theta) + s_\pi(\theta)$ of \eqref{eq:post-score-def} around the true parameter $\theta_0$:
\begin{equation}
S_{\mathrm{post}}(\hat\theta_{\mathrm{MAP}})
\;\approx\;
S_{\mathrm{post}}(\theta_0)
\;+\;
\nabla_\theta S_{\mathrm{post}}(\theta_0)\,(\hat\theta_{\mathrm{MAP}} - \theta_0).
\label{eq:posterior-score-taylor}
\end{equation}
At the MAP, $S_{\mathrm{post}}(\hat\theta_{\mathrm{MAP}}) = 0$. Replacing the empirical likelihood Hessian by its expectation $-\mathbf H_{\mathrm{lik}}(\theta_0)$ and noting that $\nabla_\theta s_\pi = -\mathbf H_{\mathrm{prior}}$ is deterministic by the convention of Section~1 (so equals its own expectation), we obtain $\nabla_\theta S_{\mathrm{post}}(\theta_0) \approx -\mathbf H_{\mathrm{post}}(\theta_0)$ and
\begin{equation}
\hat\theta_{\mathrm{MAP}} - \theta_0
\;\approx\;
\mathbf H_{\mathrm{post}}^{-1}\, S_{\mathrm{post}}(\theta_0).
\label{eq:posterior-displacement}
\end{equation}
This is the Bayesian analogue of $(\hat\theta - \theta_0) \approx \mathbf H^{-1} S(\theta_0)$ in Section~6 of the Likelihood supplement, with $\mathbf H_{\mathrm{post}}$ in place of $\mathbf H_{\mathrm{lik}}$ and the prior score $s_\pi(\theta_0)$ folded additively into $S_{\mathrm{post}}(\theta_0)$. The expansion assumes $\|\hat\theta_{\mathrm{MAP}} - \theta_0\|$ small; in the strong-prior regime examined below this displacement contracts toward $\mu_{\mathrm{prior}}$ rather than $\theta_0$, so beyond that regime $\Sigma_{\mathrm{DCC},\mathrm{post}}$ functions as a defined diagnostic (Section~6.2) rather than as the literal MAP-sampling covariance.

\subsection*{6.2 Sandwich Form: Definition of $\Sigma_{\mathrm{DCC},\mathrm{post}}$}

Taking sampling covariances on both sides of \eqref{eq:posterior-displacement} and using $\mathrm{Cov}(S_{\mathrm{post}}) = \mathbf{JT}_{\mathrm{lik}}$ from \eqref{eq:post-score-cov} (since $s_\pi$ is deterministic) would give the literal MAP sampling covariance $\mathbf H_{\mathrm{post}}^{-1}\, \mathbf{JT}_{\mathrm{lik}}\, \mathbf H_{\mathrm{post}}^{-1}$. Diagnostically, however, the prior must contribute to the information budget on the same footing as the data, exactly as in the construction of $\mathbf{JT}_{\mathrm{post}}$ in Section~3.1. We therefore adopt the augmented-observation convention and define
\begin{equation}
\Sigma_{\mathrm{DCC},\mathrm{post}}
\;:=\;
\mathbf H_{\mathrm{post}}^{-1}\,
\mathbf{JT}_{\mathrm{post}}\,
\mathbf H_{\mathrm{post}}^{-1},
\label{eq:posterior-sigma-dcc-def}
\end{equation}
the Posterior Distortion-Corrected Covariance. With $\mathbf{JT}_{\mathrm{post}} = \mathbf{JT}_{\mathrm{lik}} + \mathbf H_{\mathrm{prior}}$ from \eqref{eq:JT-post-def}, this is the direct counterpart of the Likelihood sandwich $\Sigma_{\mathrm{DCC}} = \mathbf H^{-1}\, \mathbf J_{\mathrm{total}}\, \mathbf H^{-1}$ of Section~6 of the Likelihood supplement, with both ends of the sandwich and the meat updated to their posterior counterparts.

\paragraph{Defined diagnostic versus literal sampling covariance.}
$\Sigma_{\mathrm{DCC},\mathrm{post}}$ is a defined diagnostic, not the literal MAP sampling covariance. The literal covariance derived from \eqref{eq:posterior-displacement} is $\mathbf H_{\mathrm{post}}^{-1}\, \mathbf{JT}_{\mathrm{lik}}\, \mathbf H_{\mathrm{post}}^{-1}$ and differs from \eqref{eq:posterior-sigma-dcc-def} by an additive prior-shrinkage term:
\begin{equation}
\Sigma_{\mathrm{DCC},\mathrm{post}}
\;-\;
\mathbf H_{\mathrm{post}}^{-1}\, \mathbf{JT}_{\mathrm{lik}}\, \mathbf H_{\mathrm{post}}^{-1}
\;=\;
\mathbf H_{\mathrm{post}}^{-1}\, \mathbf H_{\mathrm{prior}}\, \mathbf H_{\mathrm{post}}^{-1}.
\label{eq:posterior-sigma-dcc-gap}
\end{equation}
The two coincide under correct specification (where both reduce to $\mathbf H_{\mathrm{post}}^{-1}$, Section~6.4) and in the flat-prior limit. The augmented-observation convention is what makes the consistency reduction \eqref{eq:posterior-sigma-dcc-laplace} below hold for any prior strength.

\subsection*{6.3 Whitening Identity}

Substituting the definition \eqref{eq:posterior-c-def} of $\mathbf C_{\mathrm{post}}$ into \eqref{eq:posterior-sigma-dcc-def} yields the same algebraic identity as in the likelihood case,
\begin{equation}
\Sigma_{\mathrm{DCC},\mathrm{post}}
\;=\;
\mathbf H_{\mathrm{post}}^{-1/2}\,
\mathbf C_{\mathrm{post}}\,
\mathbf H_{\mathrm{post}}^{-1/2},
\label{eq:posterior-sigma-dcc-whitening}
\end{equation}
the Posterior analogue of $\Sigma_{\mathrm{DCC}} = \mathbf H^{-1/2}\, \mathbf C\, \mathbf H^{-1/2}$. The Posterior DCC is therefore the symmetric inverse whitening of the PID by $\mathbf H_{\mathrm{post}}^{1/2}$: $\mathbf C_{\mathrm{post}}$ encodes the directional distortion and $\mathbf H_{\mathrm{post}}^{-1/2}$ rescales it back into parameter units.

\subsection*{6.4 Consistency Reduction: Laplace Posterior Covariance}

Under the population information identity $\mathbf{JT}_{\mathrm{lik}}(\theta_0) = \mathbf H_{\mathrm{lik}}(\theta_0)$ at the true parameter (Section~4.1), $\mathbf C_{\mathrm{post}} = \mathbf I$ by \eqref{eq:posterior-c-identity}, and \eqref{eq:posterior-sigma-dcc-whitening} collapses to
\begin{equation}
\Sigma_{\mathrm{DCC},\mathrm{post}}
\;=\;
\mathbf H_{\mathrm{post}}^{-1},
\label{eq:posterior-sigma-dcc-laplace}
\end{equation}
the standard Laplace-approximation posterior covariance at the MAP. The Posterior DCC therefore reduces to the canonical Bayesian uncertainty quantification whenever the population likelihood-side information identity holds; departures from $\mathbf H_{\mathrm{post}}^{-1}$ measure precisely the posterior-relevant share of likelihood miscalibration (Section~4.2).

\subsection*{6.5 Limiting Regimes}

\paragraph{Flat-prior limit.}
As $\mathbf H_{\mathrm{prior}} \to \mathbf 0$, the substitutions \eqref{eq:flat-prior-limit} apply termwise to \eqref{eq:posterior-sigma-dcc-def}, giving
\begin{equation}
\Sigma_{\mathrm{DCC},\mathrm{post}}
\;\longrightarrow\;
\mathbf H_{\mathrm{lik}}^{-1}\,
\mathbf{JT}_{\mathrm{lik}}\,
\mathbf H_{\mathrm{lik}}^{-1}
\;=\;
\Sigma_{\mathrm{DCC}},
\label{eq:posterior-sigma-dcc-flat-limit}
\end{equation}
the Likelihood sandwich of Section~6 of the Likelihood supplement. When $\mathbf H_{\mathrm{lik}}$ is full rank this limit is literal. When $\mathbf H_{\mathrm{lik}}$ is rank-deficient, $\Sigma_{\mathrm{DCC},\mathrm{post}}$ diverges along null directions of $\mathbf H_{\mathrm{lik}}$ as $O(\lambda_{\min}(\mathbf H_{\mathrm{prior}})^{-1})$; on the retained eigenspace of $\mathbf H_{\mathrm{lik}}$ the limit recovers $\Sigma_{\mathrm{DCC}}$, and the Posterior DCC with a strictly positive $\mathbf H_{\mathrm{prior}}$ provides the regularized analogue of the subspace construction of Section~8 of the Likelihood supplement.

\paragraph{Strong-prior limit.}
In the opposite regime $\lambda_{\min}(\mathbf H_{\mathrm{prior}}) \to \infty$, the residual bound \eqref{eq:posterior-c-residual-bound} forces $\mathbf C_{\mathrm{post}} \to \mathbf I$, so by \eqref{eq:posterior-sigma-dcc-whitening} $\Sigma_{\mathrm{DCC},\mathrm{post}} \to \mathbf H_{\mathrm{post}}^{-1}$. A direct Neumann expansion of \eqref{eq:posterior-sigma-dcc-def} sharpens this to
\begin{equation}
\Sigma_{\mathrm{DCC},\mathrm{post}}
\;=\;
\mathbf H_{\mathrm{prior}}^{-1}
\;+\;
\mathbf H_{\mathrm{prior}}^{-1}\bigl(\mathbf{JT}_{\mathrm{lik}} - 2\mathbf H_{\mathrm{lik}}\bigr)\mathbf H_{\mathrm{prior}}^{-1}
\;+\;
O\!\bigl(\lambda_{\min}(\mathbf H_{\mathrm{prior}})^{-3}\bigr),
\label{eq:posterior-sigma-dcc-strong-limit}
\end{equation}
so the Posterior DCC contracts to $\mathbf H_{\mathrm{prior}}^{-1}$ at rate $O(\lambda_{\min}(\mathbf H_{\mathrm{prior}})^{-2})$ with a signed correction set by $\mathbf{JT}_{\mathrm{lik}} - 2\mathbf H_{\mathrm{lik}}$. Note that the deviation from the Laplace covariance is the prior-rescaled likelihood mismatch, $\Sigma_{\mathrm{DCC},\mathrm{post}} - \mathbf H_{\mathrm{post}}^{-1} = \mathbf H_{\mathrm{prior}}^{-1}(\mathbf{JT}_{\mathrm{lik}} - \mathbf H_{\mathrm{lik}})\mathbf H_{\mathrm{prior}}^{-1} + O(\lambda_{\min}(\mathbf H_{\mathrm{prior}})^{-3})$, vanishing at the same absolute rate as the Laplace covariance itself; the diagnostic loses absolute sensitivity to likelihood miscalibration in this regime.

\paragraph{Large-data (data-dominated) limit.}
For $N$ independent trajectories at fixed proper $\mathbf H_{\mathrm{prior}}$, $\mathbf H_{\mathrm{lik}}, \mathbf{JT}_{\mathrm{lik}} = O(N)$ while $\mathbf H_{\mathrm{prior}} = O(1)$, so
\begin{equation}
\Sigma_{\mathrm{DCC},\mathrm{post}}
\;=\;
\mathbf H_{\mathrm{lik}}^{-1}\, \mathbf{JT}_{\mathrm{lik}}\, \mathbf H_{\mathrm{lik}}^{-1}
\;+\;
O(N^{-2})
\;=\;
\Sigma_{\mathrm{DCC}} + O(N^{-2}),
\label{eq:posterior-sigma-dcc-large-n}
\end{equation}
i.e.\ $\Sigma_{\mathrm{DCC},\mathrm{post}} \propto 1/N$, recovering the Likelihood DCC scaling of Section~7 of the Likelihood supplement (Bernstein--von Mises). The same conclusion transfers to the single-trajectory $T \to \infty$ limit under stationarity and ergodicity of the per-interval scores.

\subsection*{6.6 Well-Posedness}

The Posterior DCC is well-defined whenever $\mathbf H_{\mathrm{post}} \succ \mathbf 0$. For analytic Gaussian Fisher $\mathbf H_{\mathrm{lik}} \succeq 0$ and any proper Gaussian prior, this is automatic by \eqref{eq:posterior-hessian-pd}: no subspace projection is needed, in contrast with Section~8 of the Likelihood supplement. For finite-difference $\mathbf H_{\mathrm{lik}}$ with possibly negative eigenvalues, well-posedness reduces to the dominance condition \eqref{eq:dominance-condition} of Section~2.1. In both cases $\Sigma_{\mathrm{DCC},\mathrm{post}}$ is computed via the spectral whitening of \eqref{eq:posterior-symm-sqrt}, applied symmetrically on both sides of $\mathbf{JT}_{\mathrm{post}}$; explicit inversion of $\mathbf H_{\mathrm{post}}$ is never required.

\subsection*{6.7 Role in the Bayesian Evidence Correction}

The Posterior DCC defined here is the covariance matrix that enters the distortion correction to the Bayesian model evidence developed in the companion supplement \texttt{supplement\_posterior\_evidence.tex}. Specifically, $\Sigma_{\mathrm{DCC},\mathrm{post}}$ replaces $\mathbf H_{\mathrm{post}}^{-1}$ as the local posterior covariance entering the Laplace approximation determinant for $\log p(D)$, restoring calibration under likelihood miscalibration in the same way that $\Sigma_{\mathrm{DCC}}$ restores calibration of frequentist confidence regions in the Likelihood supplement. The explicit evidence-level formula and its consequences for Bayes factor comparison are deferred to that companion document; this section provides only the parameter-space object $\Sigma_{\mathrm{DCC},\mathrm{post}}$ on which the evidence-level correction is built.

\section*{7. Posterior Distortion-Induced Bias (DIB)}

The bias counterpart of the posterior corrected covariance of Section~6 above quantifies the systematic displacement of the MAP estimator induced jointly by likelihood misspecification and prior pull. It is the direct Bayesian analogue of the Likelihood Distortion-Induced Bias of Section~7 of the Likelihood supplement, augmented by an additive prior-pull term that the likelihood-only construction cannot see.

\subsection*{7.1 Expected Posterior Score}

Taking expectations in \eqref{eq:post-score-def}--\eqref{eq:prior-score-gaussian} under the data-sampling distribution at fixed $\theta$ gives
\begin{equation}
\mathbf g_{\mathrm{post}}(\theta)
\;:=\;
\mathbb{E}[S_{\mathrm{post}}(\theta)]
\;=\;
\mathbb{E}[S_{\mathrm{lik}}(\theta)] + s_\pi(\theta)
\;=\;
\mathbf g_{\mathrm{lik}}(\theta)
- \mathbf H_{\mathrm{prior}}\bigl(\theta - \mu_{\mathrm{prior}}\bigr),
\label{eq:posterior-gpost-def}
\end{equation}
where $\mathbf g_{\mathrm{lik}}(\theta) := \mathbb{E}[S_{\mathrm{lik}}(\theta)]$ is the Likelihood DIB precursor (the $\mathbf g$ of Section~7 of the Likelihood supplement, under the bridge \eqref{eq:notation-bridge} extended to score means in the obvious way). The decomposition \eqref{eq:posterior-gpost-def} is exact and separates the data-side score drift from the prior-pull term linearly.

\paragraph{Behavior at the true parameter.}
At the data-generating $\theta_0$, the data score has zero expectation under correct likelihood specification and standard regularity,
\begin{equation}
\mathbf g_{\mathrm{lik}}(\theta_0) = \mathbf 0,
\label{eq:posterior-glik-true-zero}
\end{equation}
so $\mathbf g_{\mathrm{post}}(\theta_0) = s_\pi(\theta_0) = -\mathbf H_{\mathrm{prior}}(\theta_0 - \mu_{\mathrm{prior}})$ is generally nonzero: a correctly specified likelihood does not eliminate the prior pull, which vanishes only when the prior mean coincides with the truth ($\mu_{\mathrm{prior}} = \theta_0$). Under likelihood misspecification, $\mathbf g_{\mathrm{lik}}(\theta_0) \neq \mathbf 0$ and both contributions are present.

\subsection*{7.2 Posterior Distortion-Induced Bias}

Apply the first-order MAP expansion of Section~6.1 above to the posterior score, using the realized expansion $\hat\theta_{\mathrm{MAP}} - \theta_0 \approx \mathbf H_{\mathrm{post}}^{-1}\, S_{\mathrm{post}}(\theta_0)$ of \eqref{eq:posterior-displacement}. Taking expectations under the data-sampling distribution at $\theta_0$,
\begin{equation}
\mathbb{E}\!\left[\hat\theta_{\mathrm{MAP}} - \theta_0\right]
\;\approx\;
\mathbf H_{\mathrm{post}}^{-1}\, \mathbf g_{\mathrm{post}}(\theta_0).
\label{eq:posterior-bdib-derivation}
\end{equation}
We accordingly define the posterior distortion-induced bias
\begin{equation}
\mathbf b_{\mathrm{DIB},\mathrm{post}}(\theta)
\;:=\;
\mathbf H_{\mathrm{post}}^{-1}\, \mathbf g_{\mathrm{post}}(\theta)
\;=\;
\mathbf H_{\mathrm{post}}^{-1}\bigl(\mathbf g_{\mathrm{lik}}(\theta) + s_\pi(\theta)\bigr).
\label{eq:posterior-bdib-def}
\end{equation}
The two summands isolate the two physical mechanisms: $\mathbf H_{\mathrm{post}}^{-1}\,\mathbf g_{\mathrm{lik}}(\theta)$ is the likelihood-misspecification contribution (a posterior-rescaled version of the Likelihood DIB), and $\mathbf H_{\mathrm{post}}^{-1}\, s_\pi(\theta)$ is the prior-pull contribution, present even under correct likelihood specification whenever $\mu_{\mathrm{prior}} \neq \theta_0$. Since $\theta_0$ is unknown, the reported diagnostic substitutes the plug-in estimate $\hat\theta_{\mathrm{MAP}}$ for $\theta$ in \eqref{eq:posterior-bdib-def}; under misspecification this corresponds to evaluating $\mathbf g_{\mathrm{lik}}$ at the pseudo-true value $\theta_\ast = \arg\min \mathbb{E}[-\ell(\theta)]$ to leading order.

\paragraph{Well-posedness.}
For any proper Gaussian prior, $\mathbf H_{\mathrm{post}}^{-1}$ is unconditionally defined (Section~2.1, Eq.~\eqref{eq:posterior-hessian-pd}); for finite-difference $\mathbf H_{\mathrm{lik}}$ it is defined whenever the dominance condition \eqref{eq:dominance-condition} holds. Thus $\mathbf b_{\mathrm{DIB},\mathrm{post}}$ requires no subspace projection.

\subsection*{7.3 Sample-Size Scaling}

For $N$ independent trajectories of length $T$, the likelihood-side quantities scale with $N$ while the prior is fixed:
\begin{equation}
\mathbf H_{\mathrm{lik}} = O(N),
\qquad
\mathbf{JT}_{\mathrm{lik}} = O(N),
\qquad
\mathbf H_{\mathrm{prior}} = O(1),
\qquad
s_\pi(\theta) = O(1).
\label{eq:posterior-N-scaling-inputs}
\end{equation}
Under correct likelihood specification at $\theta_0$, $\|S_{\mathrm{lik}}(\theta_0)\| = O_p(\sqrt{N})$ while $\mathbb{E}[S_{\mathrm{lik}}(\theta_0)] = \mathbf 0$; under misspecification, persistent drift gives $\|\mathbf g_{\mathrm{lik}}(\theta_0)\| = O(N)$. Substituting into \eqref{eq:posterior-bdib-def} and using $\mathbf H_{\mathrm{post}} = \mathbf H_{\mathrm{lik}} + \mathbf H_{\mathrm{prior}} = O(N)$ in directions of positive $\mathbf H_{\mathrm{lik}}$ curvature,
\begin{equation}
\mathbf b_{\mathrm{DIB},\mathrm{post}}(\theta_0)
\;=\;
\underbrace{\mathbf H_{\mathrm{post}}^{-1}\,\mathbf g_{\mathrm{lik}}(\theta_0)}_{O(1)\ \text{(misspec)}}
\;+\;
\underbrace{\mathbf H_{\mathrm{post}}^{-1}\, s_\pi(\theta_0)}_{O(1/N)}.
\label{eq:posterior-bdib-scaling}
\end{equation}
That is, the prior-pull component is suppressed as $1/N$, while the likelihood-misspecification component approaches the Likelihood DIB asymptote $\mathbf H_{\mathrm{lik}}^{-1}\,\mathbf g_{\mathrm{lik}}$ with an $O(1/N)$ correction from $\mathbf H_{\mathrm{prior}}$. When $\mathbf H_{\mathrm{lik}}$ is rank-deficient, this asymptote is on its retained eigenspace; on the null directions of $\mathbf H_{\mathrm{lik}}$ the prior precision continues to dominate $\mathbf H_{\mathrm{post}}$ at every $N$ and the corresponding components of $\mathbf b_{\mathrm{DIB},\mathrm{post}}$ remain prior-controlled rather than approaching a Likelihood DIB limit. Comparison with the corrected variance, $\Sigma_{\mathrm{DCC},\mathrm{post}} = O(1/N)$ (\eqref{eq:posterior-sigma-dcc-large-n}), gives the familiar diagnostic: bias dominates variance as $N$ grows, exactly as in Section~7 of the Likelihood supplement, with the prior-pull contribution additionally vanishing. The same scaling transfers to the single-trajectory $T \to \infty$ regime under stationarity and ergodicity of the per-interval scores.

\subsection*{7.4 Flat-Prior and Strong-Prior Limits}

\paragraph{Flat-prior limit.}
As $\mathbf H_{\mathrm{prior}} \to \mathbf 0$, $s_\pi(\theta) \to \mathbf 0$ from \eqref{eq:prior-score-gaussian} and $\mathbf H_{\mathrm{post}} \to \mathbf H_{\mathrm{lik}}$ from \eqref{eq:flat-prior-limit}, so
\begin{equation}
\mathbf b_{\mathrm{DIB},\mathrm{post}}
\;\longrightarrow\;
\mathbf H_{\mathrm{lik}}^{-1}\,\mathbf g_{\mathrm{lik}}
\;=\;
\mathbf b_{\mathrm{DIB}},
\label{eq:posterior-bdib-flat-limit}
\end{equation}
the Likelihood DIB of Section~7 of the Likelihood supplement. When $\mathbf H_{\mathrm{lik}}$ is rank-deficient and $\mathbf g_{\mathrm{lik}}$ has nonzero projection onto its null space, $\mathbf b_{\mathrm{DIB},\mathrm{post}}$ diverges as $O(\lambda_{\min}(\mathbf H_{\mathrm{prior}})^{-1})$ along those directions; on the retained eigenspace of $\mathbf H_{\mathrm{lik}}$ the limit recovers $\mathbf b_{\mathrm{DIB}}$. For strictly positive $\mathbf H_{\mathrm{prior}}$ the construction is global, as in Section~2.1.

\paragraph{Strong-prior limit.}
In the opposite regime $\lambda_{\min}(\mathbf H_{\mathrm{prior}}) \to \infty$, the likelihood-misspecification term is dominated by the prior, $\mathbf H_{\mathrm{post}}^{-1}\,\mathbf g_{\mathrm{lik}}(\theta_0) = O(\lambda_{\min}(\mathbf H_{\mathrm{prior}})^{-1}) \to \mathbf 0$, while the prior-pull term saturates: using $s_\pi(\theta_0) = -\mathbf H_{\mathrm{prior}}(\theta_0 - \mu_{\mathrm{prior}})$ and $\mathbf H_{\mathrm{post}}^{-1} \to \mathbf H_{\mathrm{prior}}^{-1}$,
\begin{equation}
\mathbf b_{\mathrm{DIB},\mathrm{post}}(\theta_0)
\;\longrightarrow\;
-(\theta_0 - \mu_{\mathrm{prior}}),
\label{eq:posterior-bdib-strong-limit}
\end{equation}
i.e.\ the bias is the full shrinkage to the prior mean, independent of likelihood miscalibration. This is the boundary statement complementary to \eqref{eq:posterior-bdib-flat-limit}: weak priors recover the pure-likelihood diagnostic, strong priors collapse the MAP onto $\mu_{\mathrm{prior}}$.

\subsection*{7.5 Interpretation}

The decomposition \eqref{eq:posterior-bdib-def} makes the two physical sources of MAP displacement separately readable from the diagnostic, with $\theta = \hat\theta_{\mathrm{MAP}}$ used as the plug-in evaluation point (Section~7.2):
\begin{itemize}
\item $\mathbf H_{\mathrm{post}}^{-1}\, s_\pi(\hat\theta_{\mathrm{MAP}}) = -\mathbf H_{\mathrm{post}}^{-1}\mathbf H_{\mathrm{prior}}(\hat\theta_{\mathrm{MAP}} - \mu_{\mathrm{prior}})$ is the prior-pull contribution, a deterministic shrinkage toward $\mu_{\mathrm{prior}}$ exactly computable from the prior hyperparameters and $\mathbf H_{\mathrm{post}}$.
\item $\mathbf H_{\mathrm{post}}^{-1}\,\mathbf g_{\mathrm{lik}}(\hat\theta_{\mathrm{MAP}})$ is the likelihood-misspecification contribution and inherits the diagnostic role of $\mathbf b_{\mathrm{DIB}}$: it vanishes under correct specification (cf.\ \eqref{eq:posterior-glik-true-zero}) and reports persistent score drift otherwise.
\end{itemize}
A nonzero $\mathbf b_{\mathrm{DIB},\mathrm{post}}$ whose prior-pull component dominates is a benign feature of Bayesian shrinkage, vanishing with $N$; a nonzero misspecification component is the structural concern that does not vanish with sample size and is reported alongside $\Sigma_{\mathrm{DCC},\mathrm{post}}$ as the bias-variance trade-off at the relevant data scale.

\section*{8. Comparison: Singular Likelihood Hessian vs. Proper Prior}

Section~8 of the Likelihood supplement develops retained-eigenspace formulas for the Likelihood Information Distortion, DCC, DIB, and Correlation Distortion because $\mathbf H_{\mathrm{lik}}$ (the $\mathbf H$ of the Likelihood supplement, under the bridge \eqref{eq:notation-bridge}) can be singular or rank-deficient: only the subspace where $\mathbf H_{\mathrm{lik}}$ has eigenvalues above a numerical cutoff carries Gaussian Fisher information, and inverse powers of $\mathbf H_{\mathrm{lik}}$ must be restricted there. In the Posterior framework with a proper Gaussian prior, this construction is unnecessary: $\mathbf H_{\mathrm{post}} = \mathbf H_{\mathrm{lik}} + \mathbf H_{\mathrm{prior}}$ is positive definite on the full parameter space (Section~2.1), so every Posterior IDM quantity of Sections~4--7 is globally well-defined. The prior precision plays the role of regularizer that the retained-subspace projection plays in the Likelihood supplement, with the additional advantage that the prior eigenvalues provide an objective regularization scale while the Likelihood retained-eigenspace cutoff is a numerical heuristic.

\subsection*{8.1 Why no projection is required}

For a proper Gaussian prior, $\mathbf H_{\mathrm{prior}} \succ \mathbf 0$. Two cases arise depending on the sign profile of $\mathbf H_{\mathrm{lik}}$:
\begin{itemize}
\item \emph{Analytic Gaussian Fisher Information.} $\mathbf H_{\mathrm{lik}} \succeq \mathbf 0$ by construction, so $\mathbf H_{\mathrm{post}} \succeq \mathbf H_{\mathrm{prior}} \succ \mathbf 0$ unconditionally (Eq.~\eqref{eq:posterior-hessian-pd}). The symmetric inverse square root $\mathbf H_{\mathrm{post}}^{-1/2}$ used to whiten $\mathbf{JT}_{\mathrm{post}}$ and $\mathbf{JS}_{\mathrm{post}}$ is unconditionally defined on all of $\mathbb{R}^d$. Rank-deficient $\mathbf H_{\mathrm{lik}}$ (null and zero-eigenvalue directions) is handled trivially: any $\mathbf H_{\mathrm{prior}} \succ \mathbf 0$ suffices.
\item \emph{Finite-difference $\mathbf H_{\mathrm{lik}}$ with indefinite spectrum.} The Jordan decomposition $\mathbf H_{\mathrm{lik}} = \mathbf H_{\mathrm{lik}}^+ - \mathbf H_{\mathrm{lik}}^-$ of \eqref{eq:hlik-jordan-decomposition} isolates the negative part. Positive definiteness of $\mathbf H_{\mathrm{post}}$ is then equivalent to the dominance condition $\mathbf H_{\mathrm{prior}} - \mathbf H_{\mathrm{lik}}^- \succ \mathbf 0$ of \eqref{eq:dominance-condition}. This single matrix inequality is the bookkeeping condition that replaces the retained-subspace treatment of Section~8 of the Likelihood supplement: rather than discarding negatively-biased directions, the prior absorbs them.
\end{itemize}
In both cases, the symmetric inverse square root \eqref{eq:posterior-symm-sqrt} of $\mathbf H_{\mathrm{post}}$ is computed by spectral decomposition on the full parameter space; the Posterior IDM, sample distortion, correlation distortion, DCC, and DIB of Sections~4--7 are then evaluated globally without subspace selection.

\subsection*{8.2 Connection between the regimes}

The flat-prior and strong-prior regimes are continuously connected; the singular-likelihood case at vanishing prior is the only non-smooth boundary.

\paragraph{Weak prior, strong data: $\lambda_{\max}(\mathbf H_{\mathrm{prior}}) \ll \lambda_{\min}(\mathbf H_{\mathrm{lik}})$, with $\mathbf H_{\mathrm{lik}}$ full rank.}
$\mathbf H_{\mathrm{post}} \approx \mathbf H_{\mathrm{lik}}$ and $\mathbf{JT}_{\mathrm{post}} \approx \mathbf{JT}_{\mathrm{lik}}$, so $\mathbf C_{\mathrm{post}} \to \mathbf C$ entrywise. The Posterior IDM recovers the Likelihood IDM to leading order in $\mathbf H_{\mathrm{prior}}$. Analogously $\mathbf C_{\mathrm{sample},\mathrm{post}} \to \mathbf C_{\mathrm{sample}}$, $\mathbf R_{\mathrm{post}} \to \mathbf R$, $\Sigma_{\mathrm{DCC},\mathrm{post}} \to \Sigma_{\mathrm{DCC}}$, $\mathbf b_{\mathrm{DIB},\mathrm{post}} \to \mathbf b_{\mathrm{DIB}}$, as established in the flat-prior limits of Sections~4.3, 5.3, 6.5, and 7.4 above.

\paragraph{Strong prior, weak data: $\lambda_{\min}(\mathbf H_{\mathrm{prior}}) \to \infty$.}
By the residual bound \eqref{eq:posterior-c-residual-bound}, $\mathbf C_{\mathrm{post}} \to \mathbf I$ provided the likelihood-side mismatch $\mathbf{JT}_{\mathrm{lik}} - \mathbf H_{\mathrm{lik}}$ remains bounded (the regular-misspecification regime; if $\|\mathbf{JT}_{\mathrm{lik}}\|$ itself blows up with the prior, the limit need not hold). Equivalently $\Sigma_{\mathrm{DCC},\mathrm{post}} \to \mathbf H_{\mathrm{prior}}^{-1}$ (Eq.~\eqref{eq:posterior-sigma-dcc-strong-limit}) and $\mathbf b_{\mathrm{DIB},\mathrm{post}} \to -(\theta_0 - \mu_{\mathrm{prior}})$ (Eq.~\eqref{eq:posterior-bdib-strong-limit}): the diagnostic correctly reports that inference along these directions is prior-dominated, with full shrinkage to the prior mean and absolute likelihood-miscalibration signal washed out.

\paragraph{Vanishing prior with singular $\mathbf H_{\mathrm{lik}}$.}
If $\mathbf H_{\mathrm{prior}} \to \mathbf 0$ while $\mathbf H_{\mathrm{lik}}$ has a non-trivial null space (or eigenvalues below the numerical cutoff), the limit is singular. On the retained eigenspace of $\mathbf H_{\mathrm{lik}}$ (the columns of $\mathbf U_H$ whose eigenvalues exceed the cutoff, in the notation of Section~8 of the Likelihood supplement), the Posterior IDM converges to the Likelihood IDM. On the complementary null space, $\mathbf H_{\mathrm{post}}^{-1/2}$ diverges as $O(\lambda_{\min}(\mathbf H_{\mathrm{prior}})^{-1/2})$; the limit of the full matrix is undefined and the appropriate object is the retained-subspace Likelihood IDM $\mathbf C_H$ of Section~8 of the Likelihood supplement. The Posterior IDM with a strictly positive $\mathbf H_{\mathrm{prior}}$ is the regularized version of this construction: it preserves the rank-deficient directions as ``prior-dominated'' rather than excising them.

\subsection*{8.3 Practical recommendation}

The two diagnostics report complementary information and are simultaneously computable whenever both are well-defined:
\begin{itemize}
\item Use the \emph{Likelihood IDM} (Sections~4--10 of the Likelihood supplement) when $\mathbf H_{\mathrm{lik}}$ is well-conditioned and a prior-independent calibration is desired. The retained-subspace machinery of Section~8 of the Likelihood supplement is unavoidable when $\mathbf H_{\mathrm{lik}}$ is singular or near-singular.
\item Use the \emph{Posterior IDM} (Sections~4--7 above) when $\mathbf H_{\mathrm{lik}}$ is singular, near-singular, or indefinite (e.g.\ finite-difference Hessian along weakly identified directions), or when the diagnostic is to reflect the actual Bayesian posterior used for inference. The dominance condition \eqref{eq:dominance-condition} is the sole bookkeeping requirement; no eigenvalue cutoff is selected.
\item In well-conditioned regimes both are reportable, and the prior-strength contrast $\mathbf C - \mathbf C_{\mathrm{post}}$ quantifies the shrinkage effect of the prior on apparent calibration: by \eqref{eq:posterior-c-residual}, $\mathbf C_{\mathrm{post}} - \mathbf I$ is the residual $\mathbf{JT}_{\mathrm{lik}} - \mathbf H_{\mathrm{lik}}$ rescaled by $\mathbf H_{\mathrm{post}}^{-1/2}$ rather than $\mathbf H_{\mathrm{lik}}^{-1/2}$. Since $\mathbf H_{\mathrm{post}} \succeq \mathbf H_{\mathrm{lik}}$ in the Loewner order, the operator-norm bound $\|\mathbf C_{\mathrm{post}} - \mathbf I\| \leq \|\mathbf C - \mathbf I\|$ holds (the per-direction comparison is not generally monotone because $\mathbf{JT}_{\mathrm{lik}} - \mathbf H_{\mathrm{lik}}$ is indefinite under misspecification).
\end{itemize}
The choice is a reporting decision, not a structural restriction: $\mathbf C$ diagnoses pure likelihood miscalibration in Gaussian Fisher geometry; $\mathbf C_{\mathrm{post}}$ diagnoses the same miscalibration weighted by the actual posterior curvature available for inference.



\section*{9. Posterior Effective Sample Size and Correlation Distortion}

The Likelihood Correlation Distortion $\mathbf R$ of Section~9 of the Likelihood supplement isolates the temporal-dependence content of the score from its within-interval magnitude by whitening $\mathbf{JT}_{\mathrm{lik}}$ with $\mathbf{JS}_{\mathrm{lik}}^{1/2}$ rather than with $\mathbf H_{\mathrm{lik}}^{1/2}$. The Bayesian analogue uses the augmented fluctuation matrices $\mathbf{JT}_{\mathrm{post}}$ and $\mathbf{JS}_{\mathrm{post}}$ of Section~3.

\subsection*{9.1 Posterior Correlation Distortion}

Define the \emph{Posterior Correlation Distortion}
\begin{equation}
\mathbf R_{\mathrm{post}}(\theta)
\;:=\;
\mathbf{JS}_{\mathrm{post}}^{-1/2}\,
\mathbf{JT}_{\mathrm{post}}\,
\mathbf{JS}_{\mathrm{post}}^{-1/2},
\label{eq:posterior-r-def}
\end{equation}
the direct counterpart of $\mathbf R = \mathbf{JS}_{\mathrm{lik}}^{-1/2}\, \mathbf{JT}_{\mathrm{lik}}\, \mathbf{JS}_{\mathrm{lik}}^{-1/2}$ (Section~9 of the Likelihood supplement, under the bridge \eqref{eq:notation-bridge}). The symmetric inverse square root is computed by spectral decomposition of $\mathbf{JS}_{\mathrm{post}}$ as in Section~4.5; $\mathbf R_{\mathrm{post}}$ is SPD as a congruence transform of the SPD matrix $\mathbf{JT}_{\mathrm{post}}$ by $\mathbf{JS}_{\mathrm{post}}^{-1/2}$. Spectral invariants of $\mathbf R_{\mathrm{post}}$ (eigenvalues, trace, determinant) and per-parameter diagonals $[\mathbf R_{\mathrm{post}}]_{ii}$ in the fixed original parameter basis are well-defined and independent of the square-root algorithm choice (symmetric vs.\ Cholesky), in the same sense as Section~4.5.

\paragraph{Well-posedness.}
For any proper Gaussian prior, $\mathbf H_{\mathrm{prior}} \succ \mathbf 0$. Since $\mathbf{JS}_{\mathrm{lik}} \succeq \mathbf 0$ as a sum of within-interval score covariances (Section~3 of the Likelihood supplement),
\begin{equation}
\mathbf{JS}_{\mathrm{post}}
\;=\;
\mathbf{JS}_{\mathrm{lik}} + \mathbf H_{\mathrm{prior}}
\;\succeq\;
\mathbf H_{\mathrm{prior}}
\;\succ\; \mathbf 0,
\label{eq:posterior-jspost-pd}
\end{equation}
so $\mathbf{JS}_{\mathrm{post}}^{-1/2}$ is unconditionally defined, in contrast with $\mathbf{JS}_{\mathrm{lik}}^{-1/2}$, which requires $\mathbf{JS}_{\mathrm{lik}} \succ \mathbf 0$ (and otherwise the retained-subspace treatment of Section~8 of the Likelihood supplement). For analytic Gaussian Fisher $\mathbf{JT}_{\mathrm{lik}} = \mathrm{Cov}(S_{\mathrm{lik}}) \succeq \mathbf 0$, so $\mathbf{JT}_{\mathrm{post}} \succeq \mathbf H_{\mathrm{prior}} \succ \mathbf 0$ and $\mathbf R_{\mathrm{post}} \succ \mathbf 0$ globally. For finite-difference estimators where $\mathbf{JT}_{\mathrm{lik}}$ may be indefinite, the analogue of the dominance condition \eqref{eq:dominance-condition}, $\mathbf H_{\mathrm{prior}} - \mathbf{JT}_{\mathrm{lik}}^{-} \succ \mathbf 0$ (with $\mathbf{JT}_{\mathrm{lik}}^{-}$ the negative part of the estimator), is required for global SPD-ness of $\mathbf R_{\mathrm{post}}$; cf.\ Section~2.1.

\subsection*{9.2 Cancellation of the prior in the cross-interval residual}

By the sample/total decomposition \eqref{eq:posterior-sample-total} of Section~3 (and its rearrangement),
\begin{equation}
\mathbf{JT}_{\mathrm{post}} - \mathbf{JS}_{\mathrm{post}}
\;=\;
\mathbf{JT}_{\mathrm{lik}} - \mathbf{JS}_{\mathrm{lik}}
\;=\;
2\sum_{i<j}\mathrm{Cov}(s_i, s_j),
\label{eq:posterior-r-residual}
\end{equation}
since the prior precision $\mathbf H_{\mathrm{prior}}$ enters $\mathbf{JT}_{\mathrm{post}}$ and $\mathbf{JS}_{\mathrm{post}}$ identically and drops out of the difference (Section~3.2). The numerator of the unwhitened residual $\mathbf{JT}_{\mathrm{post}} - \mathbf{JS}_{\mathrm{post}}$ is therefore a pure data quantity, encoding the cross-interval covariance of the per-interval data scores $s_t$. The role of the prior in $\mathbf R_{\mathrm{post}}$ is confined to the whitening transform $\mathbf{JS}_{\mathrm{post}}^{-1/2} (\cdot)\, \mathbf{JS}_{\mathrm{post}}^{-1/2}$, which rescales the residual but does not generate any cross-interval signal of its own.

\subsection*{9.3 Independent-score limit: $\mathbf R_{\mathrm{post}} = \mathbf I$}

If the data scores are temporally uncorrelated, $\mathrm{Cov}(s_i, s_j) = \mathbf 0$ for all $i \neq j$, then $\mathbf{JT}_{\mathrm{lik}} = \mathbf{JS}_{\mathrm{lik}}$, whence $\mathbf{JT}_{\mathrm{post}} = \mathbf{JS}_{\mathrm{post}}$ and
\begin{equation}
\mathbf R_{\mathrm{post}}
\;=\;
\mathbf{JS}_{\mathrm{post}}^{-1/2}\,
\mathbf{JS}_{\mathrm{post}}\,
\mathbf{JS}_{\mathrm{post}}^{-1/2}
\;=\;
\mathbf I,
\label{eq:posterior-r-independent}
\end{equation}
exactly as in the Likelihood case. The independent-score baseline is the same for both diagnostics: a proper prior does not by itself create temporal-correlation signal in $\mathbf R_{\mathrm{post}}$.

\subsection*{9.4 Correlated-score regime: prior-softened temporal signal}

When the cross-interval sum $\sum_{i<j}\mathrm{Cov}(s_i, s_j) \neq \mathbf 0$ (the residual \eqref{eq:posterior-r-residual} is nonzero), $\mathbf{JT}_{\mathrm{post}} \neq \mathbf{JS}_{\mathrm{post}}$ and $\mathbf R_{\mathrm{post}} \neq \mathbf I$. (Individual nonzero off-diagonal covariances can cancel in the sum and yield $\mathbf R_{\mathrm{post}} = \mathbf I$ even with serial dependence; it is the residual sum that controls deviation from identity.) The deviation is governed by the same data-side residual that drives $\mathbf R - \mathbf I$ in the Likelihood case, but rescaled by the augmented whitening matrix. Subtracting the identity in \eqref{eq:posterior-r-def} and using \eqref{eq:posterior-r-residual},
\begin{equation}
\mathbf R_{\mathrm{post}} - \mathbf I
\;=\;
\mathbf{JS}_{\mathrm{post}}^{-1/2}\,
\bigl(\mathbf{JT}_{\mathrm{lik}} - \mathbf{JS}_{\mathrm{lik}}\bigr)\,
\mathbf{JS}_{\mathrm{post}}^{-1/2}.
\label{eq:posterior-r-deviation}
\end{equation}
The data residual $\mathbf{JT}_{\mathrm{lik}} - \mathbf{JS}_{\mathrm{lik}} = 2\sum_{i<j}\mathrm{Cov}(s_i,s_j)$ is identical to the one in the Likelihood case; the prior contribution enters only through $\mathbf{JS}_{\mathrm{post}}^{-1/2}$. In directions where $\mathbf H_{\mathrm{prior}}$ is large relative to $\mathbf{JS}_{\mathrm{lik}}$, the whitening is dominated by the prior precision and $\mathbf R_{\mathrm{post}} - \mathbf I$ is shrunk toward zero; in directions where $\mathbf H_{\mathrm{prior}}$ is small compared to $\mathbf{JS}_{\mathrm{lik}}$, the whitening matches the Likelihood case and $\mathbf R_{\mathrm{post}} - \mathbf I \approx \mathbf R - \mathbf I$. Informative-prior directions therefore have $\mathbf R_{\mathrm{post}} - \mathbf I$ shrunk toward zero in operator norm relative to the corresponding $\mathbf R - \mathbf I$ from the Likelihood supplement.

\subsection*{9.5 Posterior Effective Sample Size}

The variance-of-sum vs.\ sum-of-variances framework of Section~9 of the Likelihood supplement does \emph{not} transfer literally: the posterior whitening matrix $\mathbf{JS}_{\mathrm{post}} = \mathbf{JS}_{\mathrm{lik}} + \mathbf H_{\mathrm{prior}}$ contains the prior precision, so neither $\mathbf{JT}_{\mathrm{post}}$ nor $\mathbf{JS}_{\mathrm{post}}$ matches a variance-of-sum or sum-of-variances of any actual data sequence. The scalar summaries
\begin{align}
\kappa_{\mathrm{tr}}^{\mathrm{post}}
&\;=\;
\frac{1}{d}\,\mathrm{tr}(\mathbf R_{\mathrm{post}})
\;=\;
\frac{1}{d}\sum_{i=1}^d \lambda_i(\mathbf R_{\mathrm{post}}),
\label{eq:posterior-kappa-tr}\\
\kappa_{\mathrm{det}}^{\mathrm{post}}
&\;=\;
\exp\!\left(\frac{1}{d}\log\det(\mathbf R_{\mathrm{post}})\right)
\;=\;
\Bigl(\prod_{i=1}^d \lambda_i(\mathbf R_{\mathrm{post}})\Bigr)^{1/d}
\label{eq:posterior-kappa-det}
\end{align}
are spectral summaries of $\mathbf R_{\mathrm{post}}$ measuring, respectively, the average and geometric-mean per-eigen-direction inflation of the augmented within-interval budget $\mathbf{JS}_{\mathrm{post}}$ by cross-interval DATA covariance $\mathbf{JT}_{\mathrm{lik}} - \mathbf{JS}_{\mathrm{lik}}$; the prior enters only through the $\mathbf{JS}_{\mathrm{post}}$ whitening, not through the residual numerator. The corresponding posterior effective-sample-size summary is
\begin{equation}
T_{\mathrm{eff}}^{\mathrm{post}}
\;=\;
T / \kappa^{\mathrm{post}},
\qquad
\kappa^{\mathrm{post}} \in \{\kappa_{\mathrm{tr}}^{\mathrm{post}}, \kappa_{\mathrm{det}}^{\mathrm{post}}\},
\label{eq:posterior-teff}
\end{equation}
chosen according to the spectral profile. $T_{\mathrm{eff}}^{\mathrm{post}}$ does not have the literal Var$(L_X)/\sum_t$Var$(X_t)$ interpretation of the Likelihood $T_{\mathrm{eff}}$; it is a defined posterior diagnostic that measures the inflation of the per-direction posterior within-interval budget by cross-interval data covariance, not the total information available for inference.

\paragraph{Isotropic temporal dependence.}
If $\mathbf R_{\mathrm{post}} \approx \kappa^{\mathrm{post}} \mathbf I$, the data-side temporal dependence (as seen through the posterior whitening) acts isotropically and a single scalar effective sample size captures the diagnostic. Anisotropy in $\mathbf R_{\mathrm{post}}$ indicates that some eigen-directions of the posterior fluctuation budget carry more cross-interval signal than others, and the scalar summaries \eqref{eq:posterior-kappa-tr}--\eqref{eq:posterior-kappa-det} should be reported alongside the spectrum, exactly as in the Likelihood case.

\subsection*{9.6 Limiting Regimes}

Throughout this subsection, $T$ (or $N$) and $\mathbf H_{\mathrm{prior}}$ are varied independently: each limit is taken with the other quantity held fixed unless stated otherwise.

\paragraph{Flat-prior limit.}
As $\mathbf H_{\mathrm{prior}} \to \mathbf 0$ at fixed $\mathbf{JS}_{\mathrm{lik}}, \mathbf{JT}_{\mathrm{lik}}$, $\mathbf{JS}_{\mathrm{post}} \to \mathbf{JS}_{\mathrm{lik}}$ and $\mathbf{JT}_{\mathrm{post}} \to \mathbf{JT}_{\mathrm{lik}}$ by \eqref{eq:flat-prior-limit}, so
\begin{equation}
\mathbf R_{\mathrm{post}} \;\longrightarrow\; \mathbf R,
\qquad
\kappa_{\mathrm{tr}}^{\mathrm{post}} \to \kappa_{\mathrm{tr}},
\qquad
\kappa_{\mathrm{det}}^{\mathrm{post}} \to \kappa_{\mathrm{det}},
\label{eq:posterior-r-flat-limit}
\end{equation}
recovering the Likelihood Correlation Distortion of Section~9 of the Likelihood supplement. When $\mathbf{JS}_{\mathrm{lik}}$ is rank-deficient the limit is on its retained eigenspace (Section~8 of the Likelihood supplement uses its own $\mathbf{JS}_{\mathrm{lik}}$-eigenspace construction $\mathbf R_{J_{\mathrm{sample}}}$, distinct from the $\mathbf H$-eigenspace projection used for $\mathbf C$); on the null space of $\mathbf{JS}_{\mathrm{lik}}$, $\mathbf R_{\mathrm{post}}$ does not converge to $\mathbf R$ and stays prior-controlled at every $\mathbf H_{\mathrm{prior}}$. A strictly positive $\mathbf H_{\mathrm{prior}}$ regularizes both subspaces simultaneously and makes $\mathbf R_{\mathrm{post}}$ globally well-defined.

\paragraph{Strong-prior limit.}
As $\lambda_{\min}(\mathbf H_{\mathrm{prior}}) \to \infty$ at fixed $T$ (and fixed $\mathbf{JT}_{\mathrm{lik}}, \mathbf{JS}_{\mathrm{lik}}$, hence fixed bounded residual $\mathbf{JT}_{\mathrm{lik}} - \mathbf{JS}_{\mathrm{lik}}$), the operator-norm bound
\begin{equation}
\|\mathbf R_{\mathrm{post}} - \mathbf I\|
\;\leq\;
\lambda_{\min}(\mathbf{JS}_{\mathrm{post}})^{-1}\,
\bigl\|\mathbf{JT}_{\mathrm{lik}} - \mathbf{JS}_{\mathrm{lik}}\bigr\|
\;=\;
O\!\bigl(\lambda_{\min}(\mathbf H_{\mathrm{prior}})^{-1}\bigr)
\label{eq:posterior-r-strong-bound}
\end{equation}
forces $\mathbf R_{\mathrm{post}} \to \mathbf I$, hence $\kappa_{\mathrm{tr}}^{\mathrm{post}}, \kappa_{\mathrm{det}}^{\mathrm{post}} \to 1$ and $T_{\mathrm{eff}}^{\mathrm{post}} \to T$. The diagnostic correctly reports that data-side temporal dependence has no posterior leverage in the prior-dominated regime, consistent with the strong-prior collapses of $\mathbf C_{\mathrm{post}}$ (Section~5.3), $\Sigma_{\mathrm{DCC},\mathrm{post}}$ (Section~6.5), and $\mathbf b_{\mathrm{DIB},\mathrm{post}}$ (Section~7.4). $T_{\mathrm{eff}}^{\mathrm{post}} = T$ here reports prior dominance (the data-side temporal-dependence channel is whitened away) rather than data informativeness.

\paragraph{Large-data limit.}
For $N$ independent trajectories at fixed proper $\mathbf H_{\mathrm{prior}}$, $\mathbf{JT}_{\mathrm{lik}}, \mathbf{JS}_{\mathrm{lik}} = O(N)$ while $\mathbf H_{\mathrm{prior}} = O(1)$, so $\mathbf{JS}_{\mathrm{post}} = \mathbf{JS}_{\mathrm{lik}}(1 + O(1/N))$ (and similarly for $\mathbf{JT}_{\mathrm{post}}$), giving $\mathbf R_{\mathrm{post}} = \mathbf R + O(1/N)$. The posterior correlation distortion thus reproduces the Likelihood correlation distortion in the data-dominated regime, again paralleling the corresponding $\mathbf C_{\mathrm{post}}, \Sigma_{\mathrm{DCC},\mathrm{post}}, \mathbf b_{\mathrm{DIB},\mathrm{post}}$ behaviours. The convergence is on the retained eigenspace of $\mathbf{JS}_{\mathrm{lik}}$; on persistent null directions of $\mathbf{JS}_{\mathrm{lik}}$ independent of $N$, $\mathbf R_{\mathrm{post}}$ remains prior-controlled at every $N$.

\subsection*{9.7 Relation to $\mathbf C_{\mathrm{post}}$ and $\mathbf C_{\mathrm{sample},\mathrm{post}}$}

Under correct likelihood specification (Section~4.1, Eq.\ \eqref{eq:posterior-c-identity}), $\mathbf{JT}_{\mathrm{lik}} = \mathbf H_{\mathrm{lik}}$ at $\theta_0$ and $\mathbf C_{\mathrm{post}} = \mathbf I$. The companion within-interval identity $\mathbf{JS}_{\mathrm{lik}} = \mathbf H_{\mathrm{lik}}$ at the population level is the per-interval Bartlett identity (a per-observation property holding conditionally at each $t$, hence summing to the within-interval Fisher identity independently of cross-interval correlation); this is a strengthening of the total Fisher identity used in Section~4.1 and is local in time. Under the within-interval identity, $\mathbf C_{\mathrm{sample},\mathrm{post}} = \mathbf I$, and:

\begin{itemize}
\item \emph{Correct specification, no temporal correlation.} Both identities hold and $\mathbf{JT}_{\mathrm{lik}} = \mathbf{JS}_{\mathrm{lik}}$, so $\mathbf R_{\mathrm{post}} = \mathbf I = \mathbf C_{\mathrm{post}} = \mathbf C_{\mathrm{sample},\mathrm{post}}$.
\item \emph{Correct specification, temporally correlated scores.} The within-interval identity gives $\mathbf{JS}_{\mathrm{lik}} = \mathbf H_{\mathrm{lik}}$, hence $\mathbf{JS}_{\mathrm{post}} = \mathbf H_{\mathrm{post}}$ and $\mathbf C_{\mathrm{sample},\mathrm{post}} = \mathbf I$; under this identity, $\mathbf R_{\mathrm{post}} = \mathbf{JS}_{\mathrm{post}}^{-1/2}\,\mathbf{JT}_{\mathrm{post}}\,\mathbf{JS}_{\mathrm{post}}^{-1/2} = \mathbf H_{\mathrm{post}}^{-1/2}\,\mathbf{JT}_{\mathrm{post}}\,\mathbf H_{\mathrm{post}}^{-1/2} = \mathbf C_{\mathrm{post}}$, i.e.\ the two whitening operations coincide.
\item \emph{Within-interval misspecification.} $\mathbf C_{\mathrm{sample},\mathrm{post}} \neq \mathbf I$ and the relation between $\mathbf C_{\mathrm{post}}$ and $\mathbf R_{\mathrm{post}}$ is mediated by $\mathbf C_{\mathrm{sample},\mathrm{post}}$ via the sample/correlation decomposition (Section~10), which holds in the restricted sense developed there (commuting / Fisher-coordinate caveat).
\end{itemize}

The three diagnostics therefore separate three distinct posterior-side concerns: $\mathbf C_{\mathrm{post}}$ measures total miscalibration, $\mathbf C_{\mathrm{sample},\mathrm{post}}$ measures within-interval mismatch, and $\mathbf R_{\mathrm{post}}$ isolates the temporal-dependence content of the data scores as seen through the posterior whitening.

\section*{10. Posterior Sample Distortion and Decomposition}

The Posterior IDM $\mathbf C_{\mathrm{post}}$ of Section~4 combines two structurally distinct sources of miscalibration: a within-interval mismatch between sample score variance and posterior curvature (the analogue of $\mathbf C_{\mathrm{sample}}$ in Section~10 of the Likelihood supplement), and a cross-interval temporal-correlation effect (the analogue of $\mathbf R$, formalized as $\mathbf R_{\mathrm{post}}$ in Section~9 above). This section promotes the within-interval object $\mathbf C_{\mathrm{sample},\mathrm{post}}$ (introduced inline in Section~4.1 and listed among the Posterior IDM diagnostics in Section~2.2) to a labeled definition, states its consistency reduction, and establishes the sample/correlation decomposition
\begin{equation}
\mathbf C_{\mathrm{post}}
\;=\;
\mathbf M\,
\mathbf R_{\mathrm{post}}\,
\mathbf M^\top
\;=\;
\mathbf C_{\mathrm{sample},\mathrm{post}}^{1/2}\,
\widetilde{\mathbf R}_{\mathrm{post}}\,
\mathbf C_{\mathrm{sample},\mathrm{post}}^{1/2},
\label{eq:posterior-c-decomposition}
\end{equation}
where $\mathbf M$ is the (non-orthogonal) basis change between the $\mathbf{JS}_{\mathrm{post}}$-whitened and $\mathbf H_{\mathrm{post}}$-whitened bases defined below, and $\widetilde{\mathbf R}_{\mathrm{post}}$ is an orthogonal conjugate of $\mathbf R_{\mathrm{post}}$ equal to $\mathbf R_{\mathrm{post}}$ when $[\mathbf H_{\mathrm{post}}, \mathbf{JS}_{\mathrm{post}}] = \mathbf 0$. This is the structural counterpart of $\mathbf C = \mathbf C_{\mathrm{sample}}^{1/2}\,\mathbf R\,\mathbf C_{\mathrm{sample}}^{1/2}$ of Section~10 of the Likelihood supplement (which carries the analogous commuting caveat), lifted to the augmented information budget of Section~3.

\subsection*{10.1 Posterior Sample Distortion}

The Posterior Sample Distortion is the within-interval analogue of $\mathbf C_{\mathrm{post}}$, obtained by replacing $\mathbf{JT}_{\mathrm{post}}$ with $\mathbf{JS}_{\mathrm{post}}$ in the whitening transform \eqref{eq:posterior-c-def}:
\begin{equation}
\mathbf C_{\mathrm{sample},\mathrm{post}}(\theta)
\;:=\;
\mathbf H_{\mathrm{post}}^{-1/2}\,
\mathbf{JS}_{\mathrm{post}}\,
\mathbf H_{\mathrm{post}}^{-1/2}.
\label{eq:posterior-c-sample-def}
\end{equation}
By the augmented-observation convention of Section~3.1, $\mathbf{JS}_{\mathrm{post}} = \mathbf{JS}_{\mathrm{lik}} + \mathbf H_{\mathrm{prior}}$ and $\mathbf H_{\mathrm{post}} = \mathbf H_{\mathrm{lik}} + \mathbf H_{\mathrm{prior}}$ carry the prior precision identically; $\mathbf C_{\mathrm{sample},\mathrm{post}}$ therefore quantifies the within-interval mismatch between data score variance and likelihood curvature. Well-posedness: $\mathbf{JS}_{\mathrm{post}} \succeq \mathbf H_{\mathrm{prior}} \succ \mathbf 0$ for any proper Gaussian prior, so $\mathbf{JS}_{\mathrm{post}}^{-1/2}$ (needed for both $\mathbf R_{\mathrm{post}}$ and the map $\mathbf M$ below) is unconditionally defined, in parallel with $\mathbf H_{\mathrm{post}}^{-1/2}$ of Section~4.4.

\paragraph{Consistency under correct within-interval specification.}
The per-interval Bartlett identity $E[s_t s_t^\top \mid y_{1:t-1}] = -E[\nabla s_t \mid y_{1:t-1}]$ holds conditionally for each $t$ (a per-observation property of the conditional likelihood), and summing gives the within-interval Fisher identity $\mathbf{JS}_{\mathrm{lik}} = \mathbf H_{\mathrm{lik}}$ in expectation independently of cross-interval correlation. This strengthens the total Fisher identity used in Section~4.1 (which fails under serial dependence by Section~3.2) and is a genuine population statement: a derived per-interval consequence of correct conditional likelihood specification, not a quoted result from the Likelihood supplement. Substituting,
\begin{equation}
\mathbf{JS}_{\mathrm{post}}
\;=\;
\mathbf H_{\mathrm{lik}} + \mathbf H_{\mathrm{prior}}
\;=\;
\mathbf H_{\mathrm{post}},
\qquad
\mathbf C_{\mathrm{sample},\mathrm{post}}
\;=\;
\mathbf I.
\label{eq:posterior-c-sample-identity}
\end{equation}
This holds whether or not $\mathrm{Cov}(s_i, s_j) = 0$ for $i \neq j$: the temporal-correlation contribution lives entirely in the difference $\mathbf{JT}_{\mathrm{post}} - \mathbf{JS}_{\mathrm{post}} = 2\sum_{i<j}\mathrm{Cov}(s_i, s_j)$ from \eqref{eq:posterior-sample-total}, hence in $\mathbf R_{\mathrm{post}}$ rather than in $\mathbf C_{\mathrm{sample},\mathrm{post}}$. The two diagnostics separate the two failure modes through the multiplicative decomposition \eqref{eq:posterior-c-decomposition} (additive only on the log-determinant scale, Section~10.5):
\begin{itemize}
\item $\mathbf C_{\mathrm{sample},\mathrm{post}} \neq \mathbf I$ signals likelihood within-interval miscalibration ($\mathbf{JS}_{\mathrm{lik}} \neq \mathbf H_{\mathrm{lik}}$), the per-interval Fisher identity failing.
\item $\mathbf R_{\mathrm{post}} \neq \mathbf I$ signals cross-interval temporal correlation (the residual sum $\sum_{i<j}\mathrm{Cov}(s_i, s_j) \neq \mathbf 0$), present under correct within-interval specification with serially dependent scores.
\end{itemize}
Correspondingly, $\mathbf C_{\mathrm{post}} = \mathbf I$ only when both fail to fire, and otherwise $\mathbf C_{\mathrm{post}}$ inherits its deviation from one or both sources jointly.

\subsection*{10.2 Decomposition Identity}

Starting from the definitions \eqref{eq:posterior-c-def} and \eqref{eq:posterior-c-sample-def}, insert $\mathbf{JS}_{\mathrm{post}}^{1/2}\,\mathbf{JS}_{\mathrm{post}}^{-1/2} = \mathbf I$ on both sides of $\mathbf{JT}_{\mathrm{post}}$:
\begin{equation}
\mathbf C_{\mathrm{post}}
\;=\;
\mathbf H_{\mathrm{post}}^{-1/2}\,
\mathbf{JS}_{\mathrm{post}}^{1/2}\,
\bigl[\mathbf{JS}_{\mathrm{post}}^{-1/2}\,
\mathbf{JT}_{\mathrm{post}}\,
\mathbf{JS}_{\mathrm{post}}^{-1/2}\bigr]\,
\mathbf{JS}_{\mathrm{post}}^{1/2}\,
\mathbf H_{\mathrm{post}}^{-1/2}.
\label{eq:posterior-c-insert}
\end{equation}
The bracketed factor is $\mathbf R_{\mathrm{post}}$ from Section~9. Setting
\begin{equation}
\mathbf M
\;:=\;
\mathbf H_{\mathrm{post}}^{-1/2}\,
\mathbf{JS}_{\mathrm{post}}^{1/2},
\label{eq:posterior-M-def}
\end{equation}
which is generally \emph{not} symmetric (it is symmetric iff $\mathbf H_{\mathrm{post}}$ and $\mathbf{JS}_{\mathrm{post}}$ commute), Eq.\ \eqref{eq:posterior-c-insert} reads
\begin{equation}
\mathbf C_{\mathrm{post}}
\;=\;
\mathbf M\,
\mathbf R_{\mathrm{post}}\,
\mathbf M^\top,
\qquad
\mathbf M\,\mathbf M^\top
\;=\;
\mathbf H_{\mathrm{post}}^{-1/2}\,
\mathbf{JS}_{\mathrm{post}}^{1/2}\,
(\mathbf{JS}_{\mathrm{post}}^{1/2})^\top\,
(\mathbf H_{\mathrm{post}}^{-1/2})^\top
\;=\;
\mathbf H_{\mathrm{post}}^{-1/2}\,
\mathbf{JS}_{\mathrm{post}}\,
\mathbf H_{\mathrm{post}}^{-1/2}
\;=\;
\mathbf C_{\mathrm{sample},\mathrm{post}},
\label{eq:posterior-c-MMT}
\end{equation}
using symmetry of $\mathbf{JS}_{\mathrm{post}}^{1/2}$ and $\mathbf H_{\mathrm{post}}^{-1/2}$ in the third equality. So $\mathbf C_{\mathrm{sample},\mathrm{post}} = \mathbf M \mathbf M^\top$: $\mathbf M$ is a non-symmetric square-root factor of $\mathbf C_{\mathrm{sample},\mathrm{post}}$, related to the symmetric PSD root by the polar decomposition $\mathbf M = \mathbf C_{\mathrm{sample},\mathrm{post}}^{1/2}\, \mathbf Q$ with $\mathbf Q$ orthogonal. The map $\mathbf M$ is the (generally non-orthogonal) basis change from the $\mathbf{JS}_{\mathrm{post}}$-whitened basis (in which $\mathbf R_{\mathrm{post}}$ is the natural object) to the $\mathbf H_{\mathrm{post}}$-whitened basis (in which $\mathbf C_{\mathrm{post}}$ lives); conjugation by $\mathbf M$ (i.e.\ $B \mapsto \mathbf M B \mathbf M^\top$) sends $\mathbf R_{\mathrm{post}}$ to $\mathbf C_{\mathrm{post}}$.

\paragraph{Sandwich form and the polar-factor caveat.}
Substituting the polar decomposition into \eqref{eq:posterior-c-MMT},
\begin{equation}
\mathbf C_{\mathrm{post}}
\;=\;
\mathbf C_{\mathrm{sample},\mathrm{post}}^{1/2}\,
\widetilde{\mathbf R}_{\mathrm{post}}\,
\mathbf C_{\mathrm{sample},\mathrm{post}}^{1/2},
\qquad
\widetilde{\mathbf R}_{\mathrm{post}}
\;:=\;
\mathbf Q\,\mathbf R_{\mathrm{post}}\,\mathbf Q^\top,
\label{eq:posterior-c-sandwich}
\end{equation}
i.e.\ the $\mathbf C_{\mathrm{sample},\mathrm{post}}^{1/2}$-sandwich form holds with $\mathbf R_{\mathrm{post}}$ replaced by its orthogonal conjugate $\widetilde{\mathbf R}_{\mathrm{post}}$. The literal sandwich identity $\mathbf C_{\mathrm{post}} = \mathbf C_{\mathrm{sample},\mathrm{post}}^{1/2}\, \mathbf R_{\mathrm{post}}\, \mathbf C_{\mathrm{sample},\mathrm{post}}^{1/2}$ holds iff $\mathbf Q = \mathbf I$, equivalently iff $\mathbf H_{\mathrm{post}}$ and $\mathbf{JS}_{\mathrm{post}}$ commute (simultaneous diagonalizability), and in particular when $\mathbf R_{\mathrm{post}} = \kappa \mathbf I$ is scalar (since $\mathbf Q \kappa \mathbf I \mathbf Q^\top = \kappa \mathbf I$). The parallel identity of Section~10 of the Likelihood supplement, stated there ``in Fisher coordinates'', carries the same commuting caveat. The matrices $\mathbf C_{\mathrm{sample},\mathrm{post}}^{1/2}\, \mathbf R_{\mathrm{post}}\, \mathbf C_{\mathrm{sample},\mathrm{post}}^{1/2}$ and $\mathbf C_{\mathrm{post}}$ are always congruent with the same inertia and the same determinant (Section~10.5), but they are not equal as matrices in general.

\subsection*{10.3 Effective Sample Size Correction}

\paragraph{Isotropic temporal effect.}
If the temporal correlation acts isotropically in the $\mathbf{JS}_{\mathrm{post}}$-whitened basis, $\mathbf R_{\mathrm{post}} = \kappa\,\mathbf I$ for some scalar $\kappa > 0$, then orthogonal conjugation by $\mathbf Q$ leaves $\mathbf R_{\mathrm{post}}$ invariant, so the polar-factor caveat of \eqref{eq:posterior-c-sandwich} drops out and
\begin{equation}
\mathbf C_{\mathrm{post}}
\;=\;
\kappa\,\mathbf C_{\mathrm{sample},\mathrm{post}},
\qquad
\Sigma_{\mathrm{DCC},\mathrm{post}}
\;=\;
\kappa\,
\mathbf H_{\mathrm{post}}^{-1/2}\,
\mathbf C_{\mathrm{sample},\mathrm{post}}\,
\mathbf H_{\mathrm{post}}^{-1/2},
\label{eq:posterior-isotropic-kappa}
\end{equation}
using \eqref{eq:posterior-sigma-dcc-whitening}. Under the additional within-interval Fisher identity ($\mathbf C_{\mathrm{sample},\mathrm{post}} = \mathbf I$ by \eqref{eq:posterior-c-sample-identity}), this collapses to
\begin{equation}
\Sigma_{\mathrm{DCC},\mathrm{post}}
\;=\;
\kappa\,\mathbf H_{\mathrm{post}}^{-1},
\label{eq:posterior-isotropic-laplace}
\end{equation}
so the Posterior DCC equals the Laplace covariance scaled by a single effective-sample-size factor $\kappa$. The familiar $N_{\mathrm{eff}} = N/\kappa$ rescaling of frequentist standard errors lifts here to a scalar inflation of the entire posterior covariance matrix. The $\kappa\mathbf I$ ansatz applies in three limits established in Sections~4--8: (i) the strong-prior regime where $\mathbf R_{\mathrm{post}} \to \mathbf I$ and $\mathbf C_{\mathrm{sample},\mathrm{post}} \to \mathbf I$ jointly (Section~5.3), with $\kappa \to 1$; (ii) the within-interval Fisher identity (population correct specification, Section~10.1) combined with isotropic serial correlation; (iii) the flat-prior limit when the Likelihood $\kappa$ is itself near-isotropic. The brackets on $\kappa$ are: $\kappa \to \kappa_{\mathrm{lik}}$ (the Likelihood ESS) as $\mathbf H_{\mathrm{prior}} \to \mathbf 0$, and $\kappa \to 1$ as $\lambda_{\min}(\mathbf H_{\mathrm{prior}}) \to \infty$.

\paragraph{Anisotropic or sample-distorted regime.}
If $\mathbf R_{\mathrm{post}}$ is anisotropic (temporal correlation direction-dependent) or $\mathbf C_{\mathrm{sample},\mathrm{post}} \neq \mathbf I$ (within-interval likelihood mismatch), neither \eqref{eq:posterior-isotropic-kappa} nor \eqref{eq:posterior-isotropic-laplace} reduces to a scalar correction: the distortion is structurally direction-dependent and the full matrix \eqref{eq:posterior-c-decomposition}, with attention to the polar-factor caveat \eqref{eq:posterior-c-sandwich}, is required for inference. The two diagnostics serve as gates: report a scalar $N_{\mathrm{eff}}$ only when $\mathbf R_{\mathrm{post}}$ is near-isotropic and $\mathbf C_{\mathrm{sample},\mathrm{post}}$ is near $\mathbf I$ in operator norm; otherwise report $\Sigma_{\mathrm{DCC},\mathrm{post}}$ directly.

\subsection*{10.4 Flat-Prior Limit}

As $\mathbf H_{\mathrm{prior}} \to \mathbf 0$, the substitutions \eqref{eq:flat-prior-limit} apply to \eqref{eq:posterior-c-sample-def}, \eqref{eq:posterior-c-MMT}, and \eqref{eq:posterior-c-decomposition}, yielding termwise
\begin{equation}
\mathbf C_{\mathrm{sample},\mathrm{post}} \to \mathbf C_{\mathrm{sample}},
\qquad
\mathbf R_{\mathrm{post}} \to \mathbf R,
\qquad
\mathbf C_{\mathrm{post}} \to \mathbf C,
\label{eq:posterior-decomp-flat}
\end{equation}
recovering the decomposition of Section~10 of the Likelihood supplement (with the same polar-factor caveat). Rank caveats apply separately to the two whitening matrices: convergence $\mathbf C_{\mathrm{sample},\mathrm{post}} \to \mathbf C_{\mathrm{sample}}$ requires the $\mathbf H_{\mathrm{lik}}$ retained-eigenspace construction (Section~8 of the Likelihood supplement); convergence $\mathbf R_{\mathrm{post}} \to \mathbf R$ requires the distinct $\mathbf{JS}_{\mathrm{lik}}$-eigenspace construction $\mathbf R_{J_{\mathrm{sample}}}$ used there. The decomposition itself converges in the flat-prior limit on the joint retained eigenspace; off the population within-interval Fisher identity, this requires the retained eigenspaces of $\mathbf H_{\mathrm{lik}}$ and $\mathbf{JS}_{\mathrm{lik}}$ to coincide. A strictly positive $\mathbf H_{\mathrm{prior}}$ regularizes both subspaces simultaneously and makes the decomposition \eqref{eq:posterior-c-decomposition} globally well-defined.

\subsection*{10.5 Strong-Prior Limit}

As $\lambda_{\min}(\mathbf H_{\mathrm{prior}}) \to \infty$ at fixed $\mathbf H_{\mathrm{lik}}, \mathbf{JS}_{\mathrm{lik}}, \mathbf{JT}_{\mathrm{lik}}$: $\mathbf H_{\mathrm{post}} \to \mathbf H_{\mathrm{prior}}$, $\mathbf{JS}_{\mathrm{post}} \to \mathbf H_{\mathrm{prior}}$, and $\mathbf{JT}_{\mathrm{post}} \to \mathbf H_{\mathrm{prior}} + (\mathbf{JT}_{\mathrm{lik}} - \mathbf H_{\mathrm{lik}})$, so $\mathbf H_{\mathrm{post}}$ and $\mathbf{JS}_{\mathrm{post}}$ become asymptotically proportional (both approach $\mathbf H_{\mathrm{prior}}$, in particular they commute in the limit) and
\begin{equation}
\mathbf C_{\mathrm{sample},\mathrm{post}} \to \mathbf I,
\qquad
\mathbf R_{\mathrm{post}} \to \mathbf I,
\qquad
\mathbf Q \to \mathbf I,
\qquad
\mathbf C_{\mathrm{post}} \to \mathbf I,
\label{eq:posterior-decomp-strong}
\end{equation}
each with rate $O(\lambda_{\min}(\mathbf H_{\mathrm{prior}})^{-1})$ from the residual bound \eqref{eq:posterior-c-residual-bound}, in agreement with Section~5.3, Section~6.5, and Section~9.6. The decomposition \eqref{eq:posterior-c-decomposition} collapses trivially to $\mathbf C_{\mathrm{post}} = \mathbf I$. Both $\mathbf C_{\mathrm{sample},\mathrm{post}}$ and $\mathbf R_{\mathrm{post}}$ lose absolute sensitivity to likelihood miscalibration and temporal correlation in this regime — a feature, not a failure mode, since inference along prior-dominated directions is governed by the prior and not by the data-side budget.

\subsection*{10.6 Evidence-Correction Determinant Split}

Taking determinants directly from $\mathbf C_{\mathrm{post}} = \mathbf H_{\mathrm{post}}^{-1} \mathbf{JT}_{\mathrm{post}}$ (in the similarity sense) and $\mathbf C_{\mathrm{sample},\mathrm{post}} = \mathbf H_{\mathrm{post}}^{-1} \mathbf{JS}_{\mathrm{post}}$ (likewise),
\begin{equation}
\det(\mathbf C_{\mathrm{post}})
\;=\;
\frac{\det(\mathbf{JT}_{\mathrm{post}})}{\det(\mathbf H_{\mathrm{post}})}
\;=\;
\frac{\det(\mathbf{JS}_{\mathrm{post}})}{\det(\mathbf H_{\mathrm{post}})}
\cdot
\frac{\det(\mathbf{JT}_{\mathrm{post}})}{\det(\mathbf{JS}_{\mathrm{post}})}
\;=\;
\det(\mathbf C_{\mathrm{sample},\mathrm{post}})\,
\det(\mathbf R_{\mathrm{post}}),
\label{eq:posterior-det-split}
\end{equation}
which is unconditional (it does not rely on the polar-factor caveat of \eqref{eq:posterior-c-sandwich}: $\det(\mathbf M \mathbf R_{\mathrm{post}} \mathbf M^\top) = \det(\mathbf M)^2 \det(\mathbf R_{\mathrm{post}}) = \det(\mathbf C_{\mathrm{sample},\mathrm{post}}) \det(\mathbf R_{\mathrm{post}})$). This is forward-relevant for the Bayesian evidence correction developed in \texttt{supplement\_posterior\_evidence.tex}: the IDM evidence correction $\tfrac{1}{2}\log\det(\mathbf C_{\mathrm{post}})$ splits additively into a within-interval mismatch contribution $\tfrac{1}{2}\log\det(\mathbf C_{\mathrm{sample},\mathrm{post}})$ and a temporal-correlation contribution $\tfrac{1}{2}\log\det(\mathbf R_{\mathrm{post}})$. Both terms are computable separately from the same posterior IDM machinery; either can be reported as a standalone diagnostic of the corresponding failure mode, and the two are additive (not multiplicative) on the log-evidence scale. Under correct specification with no temporal correlation, both vanish; under correct within-interval specification with serial dependence only, the first vanishes and the second is the pure effective-sample-size correction; under within-interval misspecification only, the second vanishes and the first reports the per-interval Fisher mismatch.

\paragraph{Limiting behavior of the split.}
In the flat-prior limit, $\log\det(\mathbf C_{\mathrm{sample},\mathrm{post}}) \to \log\det(\mathbf C_{\mathrm{sample}})$ and $\log\det(\mathbf R_{\mathrm{post}}) \to \log\det(\mathbf R)$, recovering the Likelihood-supplement evidence split. In the strong-prior limit, both terms vanish at rate $O(\lambda_{\min}(\mathbf H_{\mathrm{prior}})^{-2})$ by \eqref{eq:posterior-decomp-strong}, in agreement with the strong-prior collapse of the absolute miscalibration signal documented in Section~6.5. The standalone-diagnostic recommendation in the previous paragraph is implicitly conditioned on the moderate-prior regime where each term is at its natural $O(1)$ scale.

\section*{11. Summary}

This supplement lifts the Likelihood Information Distortion machinery of \texttt{supplement\_information\_distortion\_main.tex} to the Bayesian posterior by treating the prior as a single deterministic augmented observation (Section~1) and propagating its precision $\mathbf H_{\mathrm{prior}}$ additively into both the curvature and the score-fluctuation budget (Sections~2--3). The resulting Posterior IDM family is globally well-defined for any proper Gaussian prior when $\mathbf H_{\mathrm{lik}}$ is PSD (analytic Gaussian Fisher), and under the dominance condition \eqref{eq:dominance-condition} when $\mathbf H_{\mathrm{lik}}$ is indefinite (finite-difference case), in both regimes without subspace projection (Section~8). It isolates pure likelihood-side miscalibration weighted by the actual posterior curvature available for inference.

\subsection*{11.1 Glossary}

\begin{align*}
\mathbf H_{\mathrm{lik}}
&\rightarrow
\text{likelihood Gaussian Fisher curvature (Section~2);}\\
&\phantom{\rightarrow}\ \text{bridges to } \mathbf H \text{ of the Likelihood supplement.}\\
\mathbf H_{\mathrm{prior}}
&\rightarrow
\text{prior precision (Section~2);}\\
&\phantom{\rightarrow}\ \text{deterministic, diagonal in transformed coordinates, absent from Likelihood IDM.}\\
\mathbf H_{\mathrm{post}}
&\rightarrow
\mathbf H_{\mathrm{lik}} + \mathbf H_{\mathrm{prior}},\ \text{posterior Hessian (Section~2);}\\
&\phantom{\rightarrow}\ \text{globally positive definite for proper prior or under the dominance condition \eqref{eq:dominance-condition}.}\\
\mathbf{JT}_{\mathrm{lik}}
&\rightarrow
\text{likelihood total score covariance (Section~3);}\\
&\phantom{\rightarrow}\ \text{bridges to } \mathbf J_{\mathrm{total}} \text{ of the Likelihood supplement.}\\
\mathbf{JS}_{\mathrm{lik}}
&\rightarrow
\text{likelihood within-interval score covariance (Section~3);}\\
&\phantom{\rightarrow}\ \text{bridges to } \mathbf J_{\mathrm{sample}} \text{ of the Likelihood supplement.}\\
\mathbf{JT}_{\mathrm{post}}
&\rightarrow
\mathbf{JT}_{\mathrm{lik}} + \mathbf H_{\mathrm{prior}}\ \text{(Section~3, augmented-observation convention);}\\
&\phantom{\rightarrow}\ \text{not the literal } \mathrm{Cov}(S_{\mathrm{post}})\text{, which is } \mathbf{JT}_{\mathrm{lik}}.\\
\mathbf{JS}_{\mathrm{post}}
&\rightarrow
\mathbf{JS}_{\mathrm{lik}} + \mathbf H_{\mathrm{prior}}\ \text{(Section~3);}\\
&\phantom{\rightarrow}\ \text{not the literal within-interval covariance of } S_{\mathrm{post}}\text{ (that is } \mathbf{JS}_{\mathrm{lik}}\text{);}\\
&\phantom{\rightarrow}\ \text{same prior augmentation as } \mathbf{JT}_{\mathrm{post}}\text{, drops out of the difference } \mathbf{JT}_{\mathrm{post}} - \mathbf{JS}_{\mathrm{post}}.\\
\mathbf C_{\mathrm{post}}
&\rightarrow
\mathbf H_{\mathrm{post}}^{-1/2}\, \mathbf{JT}_{\mathrm{post}}\, \mathbf H_{\mathrm{post}}^{-1/2},\ \text{Posterior Information Distortion (Section~4);}\\
&\phantom{\rightarrow}\ \mathbf C_{\mathrm{post}} = \mathbf I \text{ under correct likelihood specification, any proper prior.}\\
\mathbf C_{\mathrm{sample},\mathrm{post}}
&\rightarrow
\mathbf H_{\mathrm{post}}^{-1/2}\, \mathbf{JS}_{\mathrm{post}}\, \mathbf H_{\mathrm{post}}^{-1/2},\ \text{Posterior Sample Distortion (Section~10);}\\
&\phantom{\rightarrow}\ \text{within-interval mismatch only; by construction does not see cross-interval terms}\\
&\phantom{\rightarrow}\ \text{(those are isolated in } \mathbf R_{\mathrm{post}} \text{ via } \mathbf{JT}_{\mathrm{post}} - \mathbf{JS}_{\mathrm{post}}\text{).}\\
\mathbf R_{\mathrm{post}}
&\rightarrow
\mathbf{JS}_{\mathrm{post}}^{-1/2}\, \mathbf{JT}_{\mathrm{post}}\, \mathbf{JS}_{\mathrm{post}}^{-1/2},\ \text{Posterior Correlation Distortion (Section~9);}\\
&\phantom{\rightarrow}\ \text{numerator of } \mathbf R_{\mathrm{post}} - \mathbf I \text{ depends only on data scores; prior } \mathbf H_{\mathrm{prior}}\\
&\phantom{\rightarrow}\ \text{cancels from } \mathbf{JT}_{\mathrm{post}} - \mathbf{JS}_{\mathrm{post}} \text{ but rescales magnitude via } \mathbf{JS}_{\mathrm{post}}^{-1/2}.\\
\Sigma_{\mathrm{DCC},\mathrm{post}}
&\rightarrow
\mathbf H_{\mathrm{post}}^{-1}\, \mathbf{JT}_{\mathrm{post}}\, \mathbf H_{\mathrm{post}}^{-1},\ \text{Posterior Distortion-Corrected Covariance (Section~6);}\\
&\phantom{\rightarrow}\ \text{collapses to the Laplace covariance } \mathbf H_{\mathrm{post}}^{-1} \text{ under correct specification.}\\
\mathbf b_{\mathrm{DIB},\mathrm{post}}(\theta)
&\rightarrow
\mathbf H_{\mathrm{post}}^{-1}\bigl(\mathbf g_{\mathrm{lik}}(\theta) + s_\pi(\theta)\bigr),\ \text{Posterior Distortion-Induced Bias (Section~7);}\\
&\phantom{\rightarrow}\ \text{decomposes into likelihood-misspecification and prior-pull components;}\\
&\phantom{\rightarrow}\ \text{evaluated at } \theta_0 \text{ in the population derivation, at } \widehat\theta_{\mathrm{MAP}} \text{ in the plug-in.}
\end{align*}

\subsection*{11.2 Where Posterior IDM differs from Likelihood IDM}

Each Posterior IDM quantity reduces termwise to its Likelihood IDM counterpart in the flat-prior limit (Sections~4.3, 5.3, 6.5, 7.4) and in the data-dominated limit (Sections~5.3, 6.5, 7.3); when $\mathbf H_{\mathrm{lik}}$ is rank-deficient the flat-prior reduction is on its retained eigenspace, and on null directions the Posterior IDM diverges as $O(\lambda_{\min}(\mathbf H_{\mathrm{prior}})^{-1})$ as discussed in those sections. The structural differences from the Likelihood IDM are:
\begin{itemize}
\item \emph{Global well-posedness.} For any proper Gaussian prior with PSD $\mathbf H_{\mathrm{lik}}$, $\mathbf H_{\mathrm{post}} \succ \mathbf 0$ on all of $\mathbb{R}^d$ (Section~2.1), so the retained-eigenspace projection of Section~8 of the Likelihood supplement is unnecessary. For rank-deficient analytic $\mathbf H_{\mathrm{lik}}$, the strictly positive $\mathbf H_{\mathrm{prior}}$ absorbs the null directions with no ad hoc cutoff. For finite-difference $\mathbf H_{\mathrm{lik}}$ (possibly indefinite), the dominance condition \eqref{eq:dominance-condition}, $\mathbf H_{\mathrm{prior}} - \mathbf H_{\mathrm{lik}}^{-} \succ \mathbf 0$, is the sole bookkeeping requirement that handles indefiniteness without the numerical eigenvalue cutoff of the Likelihood case (Sections~2.1, 8).
\item \emph{Prior cancels from the distortion residual.} By \eqref{eq:posterior-c-residual}, $\mathbf C_{\mathrm{post}} - \mathbf I = \mathbf H_{\mathrm{post}}^{-1/2}(\mathbf{JT}_{\mathrm{lik}} - \mathbf H_{\mathrm{lik}})\mathbf H_{\mathrm{post}}^{-1/2}$, so the numerator depends only on likelihood quantities. The same cancellation holds for the difference $\mathbf{JT}_{\mathrm{post}} - \mathbf{JS}_{\mathrm{post}} = \mathbf{JT}_{\mathrm{lik}} - \mathbf{JS}_{\mathrm{lik}}$ encoded by $\mathbf R_{\mathrm{post}}$ (Sections~3.2 and 9). The prior contribution rescales the magnitude of the deviation (through $\mathbf H_{\mathrm{post}}^{-1/2}$ or $\mathbf{JS}_{\mathrm{post}}^{-1/2}$) but cannot create distortion of its own.
\item \emph{Posterior-curvature weighting.} A strong prior shrinks $\|\mathbf C_{\mathrm{post}} - \mathbf I\|$ toward zero (Section~5.3) and contracts $\Sigma_{\mathrm{DCC},\mathrm{post}}$ to $\mathbf H_{\mathrm{prior}}^{-1}$ (Section~6.5); $\mathbf b_{\mathrm{DIB},\mathrm{post}}$ evaluated at $\theta_0$ saturates to the prior-mean shrinkage $-(\theta_0 - \mu_{\mathrm{prior}})$ (Section~7.4), while its plug-in form at $\widehat\theta_{\mathrm{MAP}}$ correspondingly contracts to zero as $\widehat\theta_{\mathrm{MAP}} \to \mu_{\mathrm{prior}}$. Both diagnostics lose absolute sensitivity to likelihood miscalibration in this regime (Sections~4.2, 6.5); this is a feature, not a failure mode, since inference along those directions is governed by the prior.
\item \emph{Bias decomposition.} $\mathbf b_{\mathrm{DIB},\mathrm{post}}$ has a structural prior-pull term $\mathbf H_{\mathrm{post}}^{-1}\, s_\pi$ (this supplement, Section~7.5) absent from $\mathbf b_{\mathrm{DIB}}$; under the $N$-scaling of Section~7.3 this term is $O(1/N)$ along directions of positive $\mathbf H_{\mathrm{lik}}$ curvature, while the likelihood-misspecification term inherits the $O(1)$ behaviour of the Likelihood DIB. Along null/weakly-identified directions of $\mathbf H_{\mathrm{lik}}$, the prior-pull term remains $O(1)$ (Section~7.3).
\end{itemize}

\subsection*{11.3 Concluding remarks}

\paragraph{Posterior IDM as natural Bayesian generalization.}
The Posterior IDM is the natural Bayesian counterpart of the Likelihood IDM: it is defined whenever $\mathbf H_{\mathrm{post}} \succ \mathbf 0$ (analytic Gaussian Fisher with proper prior, or finite-difference $\mathbf H_{\mathrm{lik}}$ satisfying the dominance condition \eqref{eq:dominance-condition}), regardless of singularity or rank deficiency of $\mathbf H_{\mathrm{lik}}$. Sections~4--7 and the sample/correlation decompositions of Sections~9--10 are stated globally on $\mathbb{R}^d$ with no subspace selection.

\paragraph{Consistency identity under any proper prior.}
Under correct likelihood specification at $\theta_0$ and any proper Gaussian prior, the population information identity $\mathbf{JT}_{\mathrm{lik}} = \mathbf H_{\mathrm{lik}}$ combined with the augmented-observation convention \eqref{eq:JT-post-def} gives $\mathbf{JT}_{\mathrm{post}} = \mathbf H_{\mathrm{post}}$ by construction (not as an independent posterior-side Bartlett identity), hence $\mathbf C_{\mathrm{post}} = \mathbf I$ and $\Sigma_{\mathrm{DCC},\mathrm{post}} = \mathbf H_{\mathrm{post}}^{-1}$ (the Laplace posterior covariance) \emph{regardless of prior strength}. The prior contribution enters $\mathbf H_{\mathrm{post}}$ and $\mathbf{JT}_{\mathrm{post}}$ symmetrically and cancels under the whitening transform.

\paragraph{Diagnostic isolates pure likelihood-side miscalibration.}
Under likelihood misspecification, the deviation $\mathbf C_{\mathrm{post}} - \mathbf I = \mathbf H_{\mathrm{post}}^{-1/2}(\mathbf{JT}_{\mathrm{lik}} - \mathbf H_{\mathrm{lik}})\mathbf H_{\mathrm{post}}^{-1/2}$ has a numerator that depends only on likelihood quantities, so the diagnostic isolates likelihood-side calibration failures weighted by the actual posterior curvature. The same cancellation applies to $\mathbf{JT}_{\mathrm{post}} - \mathbf{JS}_{\mathrm{post}}$, so the numerator of $\mathbf R_{\mathrm{post}} - \mathbf I$ depends only on the data scores; the whitening by $\mathbf{JS}_{\mathrm{post}}^{-1/2} = (\mathbf{JS}_{\mathrm{lik}} + \mathbf H_{\mathrm{prior}})^{-1/2}$ rescales the magnitude in proportion to the posterior within-interval budget but cannot create cross-interval distortion of its own.

\paragraph{Likelihood vs.\ Posterior IDM as a reporting choice.}
The choice between Likelihood IDM and Posterior IDM is a reporting decision, not a structural restriction. Reported in well-conditioned regimes, $\mathbf C$ diagnoses pure likelihood miscalibration in Gaussian Fisher geometry, while $\mathbf C_{\mathrm{post}}$ diagnoses the same miscalibration weighted by the posterior curvature actually available for inference (Section~8.3). When $\mathbf H_{\mathrm{lik}}$ is singular, near-singular, or indefinite (the latter subject to the dominance condition \eqref{eq:dominance-condition}), the Posterior IDM is the only globally well-defined report and replaces the retained-eigenspace machinery of Section~8 of the Likelihood supplement.

\paragraph{Companion documents.}
The definitions established in Sections~1--10 of this supplement underwrite three further analyses in this folder. The Posterior DCC of Section~6 enters the distortion correction to the Bayesian model evidence and the Bayes factor comparison developed in \texttt{supplement\_posterior\_evidence.tex}. The full Posterior IDM family enters the information-gain decomposition developed in \texttt{supplement\_information\_gain.tex}, separating data-driven from prior-driven contributions. The dominance condition \eqref{eq:dominance-condition}, the well-posedness regimes of Section~8, and the bootstrap-level monitoring of $\mathbf C_{\mathrm{post}}$, $\mathbf R_{\mathrm{post}}$, and $\Sigma_{\mathrm{DCC},\mathrm{post}}$ feed the safety-categorization policy of \texttt{supplement\_safety\_categorization.tex}. The Posterior IDM defined here is the parameter-space substrate on which all three companion documents build.


\end{document}
